% Anexo: Entrevistas individuais completas (PRO3474)
% Transcricao integral das 11 entrevistas semiestruturadas, sem cortes.
% Pressupoe pacote inputenc/utf8 e fontenc T1 no preambulo do documento principal.

\section{Entrevista 1}

\textbf{Danilo:} Para começar, me conta um pouco sobre como é a sua relação com compra de roupa, você gosta de comprar, é mais por necessidade, funciona como um passeio?

\textbf{Marisa:} Ai, eu gosto, viu. Não é só necessidade não. Eu trabalho com venda de doce, fico em casa boa parte do dia, então quando eu saio pra comprar roupa já é tipo um respiro. Gosto de olhar tudo, mexer nas peças, mesmo que eu não vá levar nada... antigamente eu ia mais só pra passear mesmo, hoje em dia geralmente já vou com uma necessidade, tipo repor uma calça do meu marido ou ver alguma coisa pros meus enteados.

\textbf{Danilo:} Com que frequência você compra roupa em loja física, e qual unidade ou região você mais frequenta?

\textbf{Marisa:} Olha, umas duas vezes por mês eu entro numa loja, mas comprar mesmo, levar alguma coisa, é mais tipo uma vez por mês. Eu vou na Renner do Shopping Aricanduva, aqui perto de casa, na Zona Leste. É a que fica mais na mão, porque eu não gosto de pegar muito trânsito, viu.

\textbf{Danilo:} Você costuma comprar mais sozinha ou acompanhada?

\textbf{Marisa:} Ah, quase sempre com minha filha ou com minha sobrinha. Sozinha eu vou mais pra resolver rapidinho, tipo repor uma coisa básica. Mas quando é pra escolher com calma eu gosto de ir acompanhada, elas me ajuda a decidir.

\textbf{Danilo:} Isso mudou nos últimos anos?

\textbf{Marisa:} Mudou sim. Eu ia muito mais antes, quase toda semana dava uma olhada. Hoje eu tô ajudando a cuidar do meu neto, então o tempo ficou mais curto e eu preciso ser mais objetiva quando vou, né.

\textbf{Danilo:} Pensa na última vez que você comprou uma peça de roupa numa loja física. Me conta essa história do início ao fim, o que te fez ir à loja, o que aconteceu lá dentro, até você sair.

\textbf{Marisa:} Foi semana passada. Eu fui porque a calça do meu marido rasgou e ele precisava de uma pra trabalhar ainda naquela semana... deixa eu ver, foi numa terça, acho, ou quarta, não lembro direito. Cheguei lá, a loja tava bem cheia, tinha promoção de inverno. Fiquei um tempo procurando sozinha, não achei o tamanho dele na cor que eu queria, aí chamei uma moça que trabalha lá, ela é bem atenciosa, sabe, já me atendeu outras vezes. Ela foi no estoque e trouxe em outra cor, serviu bem. Aí de passagem eu vi uma blusa bonita pra mim, acabei provando também, gostei e levei. Paguei no cartão de débito e fui embora.

\textbf{Danilo:} Nesse momento, você já sabia o que estava procurando, ou foi mais explorando?

\textbf{Marisa:} Da calça eu sabia certinho, o número, a cor que ele gosta. Mas a blusa foi mais porque eu vi passando mesmo, não tava nos meus planos não.

\textbf{Danilo:} Como foi o momento de escolher e experimentar a peça?

\textbf{Marisa:} A calça eu nem provei, é dele, eu só confio no tamanho. Já a blusa eu experimentei no provador, mas tava com fila, tive que esperar um pouco pra vagar uma cabine.

\textbf{Danilo:} E o pagamento, como aconteceu?

\textbf{Marisa:} No caixa, cartão de débito, rapidinho, mas antes de chegar lá teve uma filazinha, coisa de uns dez minutos.

\textbf{Danilo:} Teve algum momento que você hesitou ou quase desistiu?

\textbf{Marisa:} Hesitei no preço da calça, tava mais caro do que eu lembrava. Mas como ele precisava mesmo, levei assim.

\textbf{Danilo:} Alguma parte demorou mais do que você esperava?

\textbf{Marisa:} O provador. Tinha só duas cabine abertas e uma fila de umas quatro pessoas.

\textbf{Danilo:} Isso é parecido com o que normalmente acontece pra você, ou foi diferente dessa vez?

\textbf{Marisa:} É bem parecido, viu. Quase sempre que eu vou tem fila de provador, principalmente fim de semana.

\textbf{Danilo:} Quando você está procurando alguma coisa específica numa loja, como você faz para encontrar?

\textbf{Marisa:} Eu ando pelos corredor olhando, mas confesso que às vezes eu fico meio perdida, a loja é grande e muda os produto de lugar. Se eu não acho rápido, eu já chamo alguém pra me ajudar, não fico procurando muito tempo sozinha não.

\textbf{Danilo:} Já aconteceu de você não encontrar o tamanho, a cor ou o modelo que queria?

\textbf{Marisa:} Demais. Principalmente pra mim mesma, que eu uso um tamanho que sempre falta. Isso é bem frequente, viu.

\textbf{Danilo:} Você pesquisa em algum canal, site, aplicativo, redes sociais, antes de ir à loja?

\textbf{Marisa:} Não, quase nunca. Eu não sou muito boa com essas aplicativo, prefiro ir direto e ver pessoalmente.

\textbf{Danilo:} O que você faz quando isso acontece, pede ajuda, procura sozinha por mais tempo, desiste?

\textbf{Marisa:} Peço ajuda quase sempre. Se não tem ninguém disponível, aí eu deixo pra outro dia.

\textbf{Danilo:} Isso te faz sair da loja de mãos vazias com que frequência?

\textbf{Marisa:} De vez em quando acontece, não é toda vez, mas acontece.

\textbf{Danilo:} Se você soubesse, antes de sair de casa, que a loja tem o que você quer, isso mudaria alguma coisa em como você compra?

\textbf{Marisa:} Mudaria bastante, viu. Eu ia direto sem ficar rodando à toa, porque às vezes eu vou, não acho, e tenho que ir embora sem nada, e isso é um saco.

\textbf{Danilo:} Voltando à sua última compra, em algum momento um vendedor te ajudou? Como foi isso?

\textbf{Marisa:} Ajudou, sim, já contei, a moça foi buscar a calça no estoque pra mim. Se não fosse ela eu tinha saído sem levar.

\textbf{Danilo:} Tem alguma parte da compra em que você prefere não falar com ninguém e resolver sozinha?

\textbf{Marisa:} Na hora de experimentar eu prefiro ficar sozinha no provedor, não gosto que fiquem esperando do lado, me sinto pressionada.

\textbf{Danilo:} O que um vendedor resolve para você que você não resolveria sozinha?

\textbf{Marisa:} Achar tamanho que não tá na araque, e também dar uma opinião se ficou bom em mim, porque às vezes eu mesma não confio no meu julgamento não.

\textbf{Danilo:} Já teve uma situação em que a ajuda de um vendedor atrapalhou mais do que ajudou, ou o contrário, sentiu falta de um vendedor e não tinha ninguém por perto?

\textbf{Marisa:} Já teve vez de eu precisar e não achar ninguém, aí eu esperei um bom tempo parada num corredor até alguém passar. Isso me deixou incomodada, viu.

\textbf{Danilo:} Isso muda dependendo do tipo de peça ou do valor?

\textbf{Marisa:} Muda. Pra coisa mais cara eu quero mais atenção, quero ter certeza que eu tô fazendo a escolha certa.

\textbf{Danilo:} Isso mudou com o tempo, você busca mais ou menos ajuda hoje do que buscava antes?

\textbf{Marisa:} Acho que busco mais hoje. Com a idade eu confio menos na minha vista pra ver se caiu bem, aí prefiro perguntar.

\textbf{Danilo:} Como foi o momento de pagar na sua última compra? Teve fila?

\textbf{Marisa:} Teve, uns dez minuto, não foi terrível mas também não foi rápido.

\textbf{Danilo:} De modo geral, o que você acha do processo entre decidir que vai levar uma peça e efetivamente sair da loja com ela?

\textbf{Marisa:} Acho meio demorado, mas eu já acostumei que é assim mesmo. Não é uma coisa que me revolta, entende, é mais uma coisa que eu já espero que aconteça.

\textbf{Danilo:} Já desistiu de comprar algo por causa da fila ou da demora no caixa?

\textbf{Marisa:} Uma vez, no fim de ano, a fila tava enorme e eu desisti, devolvi a peça na araque e fui embora.

\textbf{Danilo:} O que acontece, na prática, entre você decidir vou levar isso e você sair da loja com a peça em mãos?

\textbf{Marisa:} Eu ando até o caixa, geralmente tem fila, fico esperando, às vezes converso com quem tá do lado, e depois pago e vou.

\textbf{Danilo:} Isso acontece com frequência ou foi uma exceção?

\textbf{Marisa:} Mais em data de fim de ano e promoção. No dia a dia é mais tranquilo.

\textbf{Danilo:} Você já usou algum tipo de autoatendimento, supermercado, farmácia, cinema, sem passar por um funcionário? Como foi?

\textbf{Marisa:} Já usei o caixa de autoatendimento do mercado, mas não gosto muito não, prefiro caixa com pessoa. Fico com medo de errar alguma coisa na máquina e ninguém perceber.

\textbf{Danilo:} Você confiaria em pagar por uma peça de roupa pelo celular, sem precisar falar com ninguém na loja?

\textbf{Marisa:} Acho que não, não muito. Tenho medo de ser cobrada duas vez ou de dar algum erro e eu não saber resolver sozinha.

\textbf{Danilo:} Como você se sente ao sair de uma loja com uma compra, já teve alguma experiência desconfortável com alarme, checagem de nota ou verificação na saída?

\textbf{Marisa:} Já sim, uma vez o alarme disparou sem motivo, foi bem sem graça, todo mundo olhando. Eu tinha pago certinho mas fiquei super encabulada.

\textbf{Danilo:} O que faria você confiar mais, ou menos, num processo assim?

\textbf{Marisa:} Confiaria mais se soubesse que tem alguém pra me ajudar caso desse algum problema. Sozinha eu fico insegura.

\textbf{Danilo:} Tem alguma coisa sobre comprar roupa em loja física que a gente não conversou e que você acha importante?

\textbf{Marisa:} Acho que o provador é uma coisa que pesa bastante, viu. Poucas cabine, fila grande, isso atrapalha demais, principalmente em dia de promoção.

\textbf{Danilo:} Se você pudesse mudar uma única coisa em como as lojas funcionam hoje, o que seria?

\textbf{Marisa:} Colocaria mais vendedora e mais provador nos dias de promoção, que é quando mais precisa.

\textbf{Danilo:} Já viu, em outra loja, outro país ou outro tipo de serviço, alguma coisa que resolvesse um problema parecido com os que você mencionou?

\textbf{Marisa:} Não muito, não. Lembro de antigamente ter caderneta de fiado em loja pequena, mas isso é outra história, não tem nada a ver com isso que você tá perguntando.

\textbf{Danilo:} Voltando no que você comentou sobre faltar vendedora e provador nos dias cheios, a Riachuelo está estudando uma etiqueta de segurança na peça que passa a se comunicar com o aplicativo da loja. Ela sabe se a peça foi paga e, quando o pagamento é confirmado pelo aplicativo, ela se destrava sozinha, sem precisar passar pelo caixa. O que chama sua atenção nisso, para o bem ou para o mal?

\textbf{Marisa:} Olha, eu fico com um pé atrás, viu. Fico pensando, e se der algum erro e a etiqueta não destravar, ou pior, destravar errado e eu ser acusada de estar levando sem pagar? Isso me preocupa mais do que me anima, sinceramente.

\textbf{Danilo:} Em que situação você usaria isso, e em que situação não usaria?

\textbf{Marisa:} Acho que só usaria se tivesse uma vendedora por perto pra confirmar que deu certo. Sozinha, sem ninguém, eu não confiaria não.

\textbf{Danilo:} Isso resolve algum dos problemas que você mencionou antes, ou resolve outra coisa?

\textbf{Marisa:} Não resolve o que eu mais reclamei, que é falta de vendedora e de provador. Isso aqui é outra coisa, é mais sobre a fila do caixa, que pra mim nem é o que mais incomoda.

\textbf{Danilo:} O que precisaria existir, ou o que teria que te garantir, para você confiar nesse processo?

\textbf{Marisa:} Precisava ter uma garantia bem clara de que não vai dar problema, e alguém pra eu recorrer caso desse.

\textbf{Danilo:} O que te preocuparia mais, a tecnologia falhar, ou perder contato com alguém que te ajuda quando precisa?

\textbf{Marisa:} Perder o contato com quem ajuda, com certeza. Pra mim isso é mais importante que a tecnologia funcionar direitinho.


\section{Entrevista 2}

\textbf{Mateus:} Para começar, me conta um pouco sobre como é a sua relação com compra de roupa, você gosta de comprar, é mais por necessidade, funciona como um passeio?

\textbf{Joyce:} É mais necessidade mesmo, sinceramente. Trabalho o dia inteiro, tenho meu filho pequeno, então quando eu vou numa loja já é porque precisa mesmo, roupa dele que ficou pequena ou uma coisa minha que rasgou.

\textbf{Mateus:} Com que frequência você compra roupa em loja física, e qual unidade ou região você mais frequenta?

\textbf{Joyce:} Umas duas, três semana eu dou uma passada, sempre na Renner do Shopping Interlagos, aqui perto, na Zona Sul. É onde eu pego meu filho na escola, então já aproveito o caminho.

\textbf{Mateus:} Você costuma comprar mais sozinha ou acompanhada?

\textbf{Joyce:} Sozinha, na maioria das vezes. Às vezes levo meu filho junto, mas contando ele não como companhia, tipo, mais como responsabilidade extra.

\textbf{Mateus:} Isso mudou nos últimos anos?

\textbf{Joyce:} Mudou bastante depois que ele nasceu. Antes eu ia com mais calma, hoje é tudo correndo, entro, resolvo e saio.

\textbf{Mateus:} Pensa na última vez que você comprou uma peça de roupa numa loja física. Me conta essa história do início ao fim, o que te fez ir à loja, o que aconteceu lá dentro, até você sair.

\textbf{Joyce:} Foi tipo duas semana atrás. Fui comprar um casaco pro meu filho porque tava esfriando e ele não tinha nada quentinho. Cheguei, fui direto na seção infantil, já sabia mais ou menos o que queria, achei rápido, provei nele ali mesmo, coloquei o casaco e vi que servia, e fui pagar.

\textbf{Mateus:} Me dá um exemplo de como foi essa parte de achar o casaco, foi tranquilo mesmo?

\textbf{Joyce:} Foi, porque eu fui num dia de semana de manhã, loja meio vazia. Achei umas três opção de cor, escolhi a mais neutra, testei nele e já resolvi.

\textbf{Mateus:} Nesse momento, você já sabia o que estava procurando, ou foi mais explorando?

\textbf{Joyce:} Sabia certinho, casaco, tamanho dele, não tinha muito tempo pra ficar olhando outra coisa.

\textbf{Mateus:} Como foi o momento de escolher e experimentar a peça?

\textbf{Joyce:} Rápido, testei nele ali mesmo no corredor, nem fui no provador porque não tinha necessidade.

\textbf{Mateus:} E o pagamento, como aconteceu?

\textbf{Joyce:} Paguei no cartão, no aplicativo do banco mesmo, por aproximação. Não teve fila porque era horário fraco.

\textbf{Mateus:} Teve algum momento que você hesitou ou quase desistiu?

\textbf{Joyce:} Não, dessa vez não. Foi tranquilo.

\textbf{Mateus:} Alguma parte demorou mais do que você esperava?

\textbf{Joyce:} Não, foi rápido, uns quinze minuto do início ao fim.

\textbf{Mateus:} Isso é parecido com o que normalmente acontece pra você, ou foi diferente dessa vez?

\textbf{Joyce:} Foi mais tranquilo que o normal, porque geralmente eu vou fim de semana, que é quando dá pra eu ter mais tempo, e aí tem mais gente.

\textbf{Mateus:} Quando você está procurando alguma coisa específica numa loja, como você faz para encontrar?

\textbf{Joyce:} Eu ando pelas seção, já tenho noção de onde fica cada coisa porque eu vou sempre na mesma loja. Se não acho eu procuro alguém rapidinho, mas prefiro tentar sozinha primeiro.

\textbf{Mateus:} Já aconteceu de você não encontrar o tamanho, a cor ou o modelo que queria?

\textbf{Joyce:} Já, principalmente roupa de criança em tamanho específico, eles vende muito rápido.

\textbf{Mateus:} Você pesquisa em algum canal, site, aplicativo, redes sociais, antes de ir à loja?

\textbf{Joyce:} Às vezes eu dou uma olhada no site antes, mas não confio cem por cento que o que tá lá vai estar na loja de verdade. Aí eu vou mesmo assim conferir pessoalmente.

\textbf{Mateus:} Conta mais sobre isso, por que você não confia que o que tá no site vai estar na loja?

\textbf{Joyce:} Porque já aconteceu de eu ver uma coisa disponível no site e chegar lá e não achar, aí acabei desconfiando, sabe. Prefiro ir olhar com meus próprios olho.

\textbf{Mateus:} O que você faz quando isso acontece, pede ajuda, procura sozinha por mais tempo, desiste?

\textbf{Joyce:} Pergunto pra alguém, se não tiver ninguém e eu tiver com pressa, que é o mais comum, eu desisto e vou embora, não tenho tempo de ficar esperando não.

\textbf{Mateus:} Isso te faz sair da loja de mãos vazias com que frequência?

\textbf{Joyce:} De vez em quando, tipo uma vez por mês isso acontece.

\textbf{Mateus:} Se você soubesse, antes de sair de casa, que a loja tem o que você quer, isso mudaria alguma coisa em como você compra?

\textbf{Joyce:} Mudaria muito, porque eu tenho pouco tempo, se eu já soubesse que tem, eu ia direto sem perder tempo procurando.

\textbf{Mateus:} Voltando à sua última compra, em algum momento um vendedor te ajudou? Como foi isso?

\textbf{Joyce:} Não precisei dessa vez, resolvi tudo sozinha.

\textbf{Mateus:} Tem alguma parte da compra em que você prefere não falar com ninguém e resolver sozinha?

\textbf{Joyce:} Praticamente tudo, na verdade. Eu gosto de resolver rápido sem precisar ficar explicando o que eu quero pra alguém.

\textbf{Mateus:} O que um vendedor resolve para você que você não resolveria sozinha?

\textbf{Joyce:} Quando não acho um tamanho na araque, aí eu preciso perguntar se tem no estoque.

\textbf{Mateus:} Já teve uma situação em que a ajuda de um vendedor atrapalhou mais do que ajudou, ou o contrário, sentiu falta de um vendedor e não tinha ninguém por perto?

\textbf{Joyce:} Já senti falta sim, uma vez eu precisava de um tamanho e não achei ninguém por dez minuto, isso me irritou porque eu tava com pressa.

\textbf{Mateus:} Isso muda dependendo do tipo de peça ou do valor?

\textbf{Joyce:} Um pouco. Pra roupa de criança eu me importo mais em achar rápido, pra roupa minha eu até gosto de olhar mais com calma quando dá tempo.

\textbf{Mateus:} Como foi o momento de pagar na sua última compra? Teve fila?

\textbf{Joyce:} Não teve fila dessa vez, mas geralmente tem, principalmente fim de semana.

\textbf{Mateus:} De modo geral, o que você acha do processo entre decidir que vai levar uma peça e efetivamente sair da loja com ela?

\textbf{Joyce:} Acho que podia ser mais rápido. Pra quem tem pouco tempo como eu, cada minuto de fila pesa.

\textbf{Mateus:} Já desistiu de comprar algo por causa da fila ou da demora no caixa?

\textbf{Joyce:} Já, mais de uma vez. Se a fila tá enorme e eu tenho hora pra buscar meu filho, eu devolvo a peça e vou embora.

\textbf{Mateus:} Me conta de uma vez específica que isso aconteceu.

\textbf{Joyce:} Foi num sábado, eu tava com uma blusa escolhida, vi a fila do caixa enorme, olhei no relógio, vi que ia atrasar, e simplesmente deixei a peça numa araque perto e saí.

\textbf{Mateus:} O que acontece, na prática, entre você decidir vou levar isso e você sair da loja com a peça em mãos?

\textbf{Joyce:} Eu ando até o caixa, olho o tamanho da fila, se for rápido eu fico, se não for eu já penso duas vez.

\textbf{Mateus:} Isso acontece com frequência ou foi uma exceção?

\textbf{Joyce:} Acontece bastante, principalmente fim de semana e data comemorativa.

\textbf{Mateus:} Você já usou algum tipo de autoatendimento, supermercado, farmácia, cinema, sem passar por um funcionário? Como foi?

\textbf{Joyce:} Uso direto no mercado, autoatendimento, acho ótimo, bem mais rápido.

\textbf{Mateus:} Você confiaria em pagar por uma peça de roupa pelo celular, sem precisar falar com ninguém na loja?

\textbf{Joyce:} Confiaria, eu já uso pix e cartão pelo celular pra tudo, não vejo muito problema nisso não.

\textbf{Mateus:} Como você se sente ao sair de uma loja com uma compra, já teve alguma experiência desconfortável com alarme, checagem de nota ou verificação na saída?

\textbf{Joyce:} Nunca tive problema com alarme, mas às vezes acho chato ter que mostrar a nota na saída, parece que estão desconfiando de mim.

\textbf{Mateus:} O que faria você confiar mais, ou menos, num processo assim?

\textbf{Joyce:} Confiaria mais se fosse rápido e não desse trabalho, tipo, se eu não precisasse ficar mostrando nada pra ninguém.

\textbf{Mateus:} Tem alguma coisa sobre comprar roupa em loja física que a gente não conversou e que você acha importante?

\textbf{Joyce:} Acho que falta mais opção de horário com loja mais vazia. Eu queria poder ir de manhã cedo sempre, mas nem sempre dá.

\textbf{Mateus:} Se você pudesse mudar uma única coisa em como as lojas funcionam hoje, o que seria?

\textbf{Joyce:} Deixaria o caixa mais rápido, com menos gente na frente.

\textbf{Mateus:} Já viu, em outra loja, outro país ou outro tipo de serviço, alguma coisa que resolvesse um problema parecido com os que você mencionou?

\textbf{Joyce:} Vi um mercado aqui perto que tem uns totem de autoatendimento bem rápido, gostaria que loja de roupa tivesse algo assim também.

\textbf{Mateus:} Você comentou que gostaria de um jeito mais rápido de pagar, sem tanta fila. A Riachuelo está estudando uma etiqueta de segurança na peça que passa a se comunicar com o aplicativo da loja. Ela sabe se a peça foi paga e, quando o pagamento é confirmado pelo aplicativo, ela se destrava sozinha, sem precisar passar pelo caixa. O que chama sua atenção nisso, para o bem ou para o mal?

\textbf{Joyce:} Achei interessante, principalmente pra quando eu tô com pressa. Se funcionar direito, resolveria bastante coisa pra mim.

\textbf{Mateus:} Em que situação você usaria isso, e em que situação não usaria?

\textbf{Joyce:} Usaria em qualquer compra rápida, tipo roupa básica. Talvez não usasse se fosse alguma coisa mais cara, aí eu ia querer ter certeza de que deu tudo certo antes de sair.

\textbf{Mateus:} Isso resolve algum dos problemas que você mencionou antes, ou resolve outra coisa?

\textbf{Joyce:} Resolve bem o problema da fila, que é o que mais me incomoda. Não resolve o problema de achar tamanho, mas pelo menos essa parte ajudaria.

\textbf{Mateus:} O que precisaria existir, ou o que teria que te garantir, para você confiar nesse processo?

\textbf{Joyce:} Precisava mostrar bem claro no aplicativo que o pagamento foi confirmado, algum aviso, tipo pode sair, tá tudo certo.

\textbf{Mateus:} Você entende isso como uma mudança só na trava da peça, ou como algo que muda a experiência de compra como um todo?

\textbf{Joyce:} Acho que muda mais a parte de sair rápido da loja. O resto, achar produto, provar, continua igual.

\textbf{Mateus:} Isso mudaria a forma como você usa o vendedor hoje?

\textbf{Joyce:} Não muito, eu já uso pouco vendedor mesmo, então continuaria mais ou menos igual.


\section{Entrevista 3}

\textbf{Lucas:} E aí, Robério, bora começar. Me conta um pouco, como é a sua relação com compra de roupa? Você gosta, é mais por necessidade, é tipo um passeio?

\textbf{Robério:} Ah, é mais necessidade mesmo. Eu não sou muito de sair pra comprar roupa por diversão não. Quando eu vou é porque precisa.

\textbf{Lucas:} Com que frequência você compra, e onde costuma ir?

\textbf{Robério:} Poucas vez, viu. Umas três, quatro vez por ano. Vou na C\&A do Shopping Cidade Sorocaba mesmo, perto de casa.

\textbf{Lucas:} E costuma ir sozinho ou acompanhado?

\textbf{Robério:} Geralmente com minha esposa. Ela que me ajuda a escolher, porque sozinho eu me perco um pouco.

\textbf{Lucas:} Isso mudou nos últimos anos?

\textbf{Robério:} Mudou. Antes eu ia mais, trabalhava, precisava de roupa pra semana toda. Aposentei, aí ficou bem mais raro.

\textbf{Lucas:} Beleza. Agora pensa na última vez que você comprou uma roupa numa loja física. Me conta essa história inteira, desde o que te fez ir até você sair de lá.

\textbf{Robério:} Foi acho que uns dois mês atrás. Minha esposa que me levou, precisava de uma camisa pra um casamento. A gente foi, ela que escolheu umas opção, eu só fui provando o que ela trazia.

\textbf{Lucas:} E você já sabia o que queria, ou foi mais ela que foi guiando?

\textbf{Robério:} Foi ela que guiou, eu não entendo muito dessas coisa de moda não.

\textbf{Lucas:} Como foi a parte de experimentar?

\textbf{Robério:} Fui no provador, ela ficava esperando do lado de fora, ia passando as opção.

\textbf{Lucas:} E o pagamento, como foi?

\textbf{Robério:} No caixa, no cartão. Teve uma filazinha, mas nada absurdo.

\textbf{Lucas:} Teve algum momento que você hesitou ou quase desistiu?

\textbf{Robério:} Hesitei no preço, achei meio caro, mas como era pra ocasião especial, levei.

\textbf{Lucas:} Alguma parte demorou mais do que esperava?

\textbf{Robério:} Não, foi tranquilo, uns quarenta minuto no total.

\textbf{Lucas:} E isso é parecido com o que normalmente rola, ou foi diferente?

\textbf{Robério:} É parecido. Eu não vou com muita frequência, mas quando vou é sempre assim, minha esposa na frente decidindo e eu ajudando pouco, viu.

\textbf{Lucas:} Quando você procura alguma coisa específica, como você faz pra achar?

\textbf{Robério:} Eu mesmo não procuro muito não, prefiro perguntar pra alguém da loja logo de cara.

\textbf{Lucas:} Já aconteceu de não achar o tamanho ou modelo que queria?

\textbf{Robério:} Já, sim. Às vezes o tamanho que eu uso já tinha acabado.

\textbf{Lucas:} Você pesquisa em site ou aplicativo antes de ir?

\textbf{Robério:} Não, isso eu não mexo. Nem tenho muita paciência pra essas coisa de celular, viu.

\textbf{Lucas:} E quando não acha o que quer, o que você faz? Pede ajuda, procura mais, desiste?

\textbf{Robério:} Peço ajuda pra minha esposa primeiro, e se ela não resolver a gente chama alguém da loja.

\textbf{Lucas:} Isso faz você sair de mãos vazias com que frequência?

\textbf{Robério:} De vez em quando acontece, mas não é o mais comum não.

\textbf{Lucas:} Se você soubesse antes que a loja tem o que você quer, isso mudaria alguma coisa?

\textbf{Robério:} Mudaria, porque eu não gosto de ir à toa. Se eu soubesse que não tem, nem ia.

\textbf{Lucas:} Voltando pra última compra, um vendedor te ajudou em algum momento?

\textbf{Robério:} Ajudou sim, minha esposa chamou uma moça pra trazer outro tamanho.

\textbf{Lucas:} Tem alguma parte que você prefere resolver sozinho, sem falar com ninguém?

\textbf{Robério:} Olha, eu prefiro que tenha sempre alguém por perto, viu. Não gosto muito de ficar resolvendo essas coisa sozinho.

\textbf{Lucas:} O que um vendedor resolve pra você que você não resolveria sozinho?

\textbf{Robério:} Quase tudo, sinceramente. Escolher, saber se ficou bom, achar tamanho. Eu confio mais neles do que em mim mesmo nessas hora.

\textbf{Lucas:} Já teve alguma vez que a ajuda atrapalhou, ou o contrário, sentiu falta e não tinha ninguém?

\textbf{Robério:} Já senti falta, uma vez fiquei um tempo esperando alguém aparecer e fiquei sem graça de ficar ali parado.

\textbf{Lucas:} Isso muda dependendo da peça ou do valor?

\textbf{Robério:} Acho que não muda muito não, eu sempre prefiro ter ajuda, seja coisa cara ou barata.

\textbf{Lucas:} E mudou com o tempo? Você busca mais ou menos ajuda hoje?

\textbf{Robério:} Acho que sempre busquei bastante, não mudou muito não.

\textbf{Lucas:} Se não tivesse nenhum vendedor na loja inteira, o que você faria diferente?

\textbf{Robério:} Acho que eu ia embora sem comprar, viu. Não ia me sentir seguro de escolher sozinho não.

\textbf{Lucas:} Como foi o pagamento na sua última compra? Teve fila?

\textbf{Robério:} Teve uma filazinha pequena, nada demorado.

\textbf{Lucas:} De modo geral, o que você acha desse processo entre decidir levar a peça e sair da loja?

\textbf{Robério:} Acho normal, não me incomoda muito não. Pra mim o que mais demora é a escolha mesmo, não o pagamento.

\textbf{Lucas:} Já desistiu de comprar algo por causa da fila?

\textbf{Robério:} Não, nunca desisti por causa disso.

\textbf{Lucas:} O que acontece, na prática, entre você decidir vou levar isso e sair da loja com a peça?

\textbf{Robério:} Vou até o caixa, espero a vez, pago e saio. É bem direto pra mim.

\textbf{Lucas:} Você já usou algum tipo de autoatendimento, tipo mercado ou farmácia, sem passar por funcionário?

\textbf{Robério:} Já vi tendo em mercado, mas eu mesmo prefiro sempre o caixa com pessoa. Não confio muito nessas máquina não.

\textbf{Lucas:} Você confiaria em pagar por uma roupa pelo celular, sem falar com ninguém na loja?

\textbf{Robério:} Não, isso eu não confiaria muito não. Prefiro ter uma pessoa por perto pra garantir que tá tudo certo.

\textbf{Lucas:} Já teve alguma experiência desconfortável com alarme ou checagem na saída?

\textbf{Robério:} Não, nunca tive problema desse tipo.

\textbf{Lucas:} O que faria você confiar mais, ou menos, num processo assim?

\textbf{Robério:} Confiaria mais se tivesse sempre alguém que eu pudesse perguntar se desse algum problema.

\textbf{Lucas:} Tem alguma coisa sobre comprar roupa que a gente não conversou e você acha importante?

\textbf{Robério:} Acho que já falamos do principal, que é eu gostar de ter gente que ajude.

\textbf{Lucas:} Se pudesse mudar uma coisa só nas lojas hoje, o que seria?

\textbf{Robério:} Colocaria mais gente pra atender, principalmente perto do provador.

\textbf{Lucas:} Já viu em outro lugar alguma coisa parecida que resolvesse algum desses problemas?

\textbf{Robério:} Não, não vi não. Eu não saio muito, nem presto muita atenção nessas coisa não.

\textbf{Lucas:} Legal. Deixa eu te falar de uma ideia que a Riachuelo tá estudando, uma etiqueta de segurança na peça que se comunica com o aplicativo da loja. Ela sabe se a peça foi paga, e quando o pagamento é confirmado pelo aplicativo, ela se destrava sozinha, sem precisar passar pelo caixa. O que você acha disso, pra bem ou pra mal?

\textbf{Robério:} Ah, isso me deixa bem inseguro, viu. Eu nem uso muito celular pra essas coisa, imagina pagar sozinho e confiar que vai destravar certo.

\textbf{Lucas:} Em que situação você usaria isso, e em que não usaria?

\textbf{Robério:} Acho que eu não usaria em quase nenhuma, sinceramente. Prefiro sempre o jeito de sempre, com pessoa ajudando.

\textbf{Lucas:} Isso resolve algum dos problemas que você comentou, ou é outra coisa?

\textbf{Robério:} Não resolve o meu problema, que é precisar de ajuda pra escolher e confiar na decisão. Isso aí é sobre pagar, e pagar nunca foi o que mais me incomodou.

\textbf{Lucas:} O que precisaria pra você confiar nesse processo?

\textbf{Robério:} Ter alguém do lado pra me ajudar a usar, pelo menos as primeira vez.

\textbf{Lucas:} O que te preocuparia mais, a tecnologia falhar, ou perder o contato com quem te ajuda?

\textbf{Robério:} Perder o contato com quem ajuda, com certeza. Isso pra mim é o mais importante.


\section{Entrevista 4}

\textbf{Mateus:} Vamos começar. Como é sua relação com compra de roupa? Gosta de comprar, é mais necessidade, ou é tipo um passeio?

\textbf{Felipe:} É bem prático pra mim. Eu não curto muito ficar rodando loja, vou quando preciso repor alguma coisa, camisa social, meia, essas coisa básica.

\textbf{Mateus:} Com que frequência você compra, e onde costuma ir?

\textbf{Felipe:} Umas seis, sete vez por ano eu acho. Vou na Zara do Shopping Eldorado, em Pinheiros, perto do trabalho.

\textbf{Mateus:} Sozinho ou acompanhado, geralmente?

\textbf{Felipe:} Sozinho, sempre. Costumo ir no horário de almoço mesmo, rapidinho.

\textbf{Mateus:} Isso mudou nos últimos anos?

\textbf{Felipe:} Um pouco. Antes eu comprava mais pelo site, mas desisti porque errava tamanho direto. Hoje prefiro ir na loja mesmo, mais certeiro.

\textbf{Mateus:} Interessante isso, me conta mais sobre esse período que você comprava pelo site e desistiu.

\textbf{Felipe:} Ah, comprei umas três vez calça e camisa, duas vez veio errado o caimento, tive que trocar. Depois disso achei mais garantido ir na loja e já sair com a certeza.

\textbf{Mateus:} Pensa na última vez que você comprou uma peça numa loja física. Me conta essa história do início ao fim.

\textbf{Felipe:} Foi semana passada. Precisava de umas camisa social pra trabalho. Entrei, fui direto na seção masculina, já sabia o tamanho, peguei três modelo, fui no provador rapidinho, três minuto, escolhi duas, e fui pagar.

\textbf{Mateus:} Você já sabia o que estava procurando, ou foi mais explorando?

\textbf{Felipe:} Sabia certinho. Não tenho paciência de ficar explorando à toa.

\textbf{Mateus:} Como foi o momento de escolher e experimentar?

\textbf{Felipe:} Rápido. Provei as três, vi qual caía melhor, decidi na hora.

\textbf{Mateus:} E o pagamento, como aconteceu?

\textbf{Felipe:} Cartão de crédito, sem fila, dia de semana de tarde é mais vazio.

\textbf{Mateus:} Teve algum momento que você hesitou ou quase desistiu?

\textbf{Felipe:} Não, foi bem direto dessa vez.

\textbf{Mateus:} Alguma parte demorou mais do que esperava?

\textbf{Felipe:} Não, foi rápido, uns quinze minuto.

\textbf{Mateus:} Isso é parecido com o normal pra você, ou foi diferente?

\textbf{Felipe:} É bem parecido, eu sempre tento ser rápido, vou com objetivo claro.

\textbf{Mateus:} Quando você procura alguma coisa específica, como você faz pra achar?

\textbf{Felipe:} Eu já sei mais ou menos onde fica cada seção, ando direto pra lá. Se não acho de primeira, procuro sozinho mais um pouco antes de pedir ajuda.

\textbf{Mateus:} Já aconteceu de não achar o tamanho ou modelo que queria?

\textbf{Felipe:} Já, às vezes acaba o tamanho P ou GG, que são os que mais saem.

\textbf{Mateus:} Você pesquisa em site ou aplicativo antes de ir?

\textbf{Felipe:} Olho às vezes, mas o app eu acho meio bagunçado pra navegar, prefiro ir direto na loja ver.

\textbf{Mateus:} Conta mais sobre isso, o que exatamente te incomoda no app?

\textbf{Felipe:} Acho a navegação confusa, demora pra carregar, às vezes eu nem sei se o filtro que eu apliquei realmente reflete o que tem na loja física perto de mim. Aí eu já nem uso mais.

\textbf{Mateus:} O que você faz quando não acha o que quer na loja, pede ajuda, procura mais, desiste?

\textbf{Felipe:} Procuro mais um pouco, e se realmente não achar eu só vou embora, não fico insistindo.

\textbf{Mateus:} Isso te faz sair de mãos vazias com que frequência?

\textbf{Felipe:} De vez em quando, talvez uma em cada cinco vez que eu vou.

\textbf{Mateus:} Se você soubesse antes que a loja tem o que você quer, isso mudaria alguma coisa?

\textbf{Felipe:} Mudaria bastante, eu economizaria tempo, que é o que eu mais valorizo nessa parada.

\textbf{Mateus:} Voltando à sua última compra, um vendedor te ajudou em algum momento?

\textbf{Felipe:} Não precisei dessa vez.

\textbf{Mateus:} Tem alguma parte que você prefere resolver sozinho?

\textbf{Felipe:} Praticamente tudo. Não gosto muito de ficar sendo abordado toda hora que eu entro numa loja.

\textbf{Mateus:} O que um vendedor resolve pra você que você não resolveria sozinho?

\textbf{Felipe:} Às vezes eu peço opinião sobre caimento, tipo se uma calça social ficou boa ou não, porque eu mesmo tenho dificuldade de julgar isso no espelho.

\textbf{Mateus:} Interessante, então mesmo preferindo resolver sozinho, tem uma parte que você recorre a alguém?

\textbf{Felipe:} É, exatamente. Pra coisa mais formal, tipo terno, calça social, eu gosto de ouvir uma segunda opinião. Pra básico, camiseta, cueca, eu resolvo tudo sozinho.

\textbf{Mateus:} Já teve uma situação em que a ajuda de um vendedor atrapalhou mais do que ajudou, ou sentiu falta e não tinha ninguém?

\textbf{Felipe:} Já tive vendedor meio insistente demais, ficando do lado sugerindo coisa que eu não pedi, isso me incomoda.

\textbf{Mateus:} Isso muda dependendo do tipo de peça ou do valor?

\textbf{Felipe:} Muda, como falei, pra roupa social eu quero mais opinião, pro resto eu já sei o que quero.

\textbf{Mateus:} Isso mudou com o tempo?

\textbf{Felipe:} Acho que sim, hoje eu sei melhor o que quero, então preciso de menos ajuda do que quando era mais novo.

\textbf{Mateus:} Como foi o pagamento na sua última compra? Teve fila?

\textbf{Felipe:} Não teve, mas normalmente tem sim, principalmente no almoço, que é quando eu costumo ir.

\textbf{Mateus:} De modo geral, o que você acha do processo entre decidir levar a peça e sair da loja?

\textbf{Felipe:} Acho que é onde mais se perde tempo. Se eu já escolhi, eu queria só pagar e sair, mas às vezes a fila estraga isso.

\textbf{Mateus:} Já desistiu de comprar algo por causa da fila ou demora no caixa?

\textbf{Felipe:} Já sim, mais de uma vez. Se eu vejo fila grande na hora do almoço eu simplesmente deixo a peça e volto outro dia, ou nem volto mais.

\textbf{Mateus:} Me conta de uma vez específica.

\textbf{Felipe:} Foi um dia que eu só tinha meia hora de almoço, escolhi uma camisa rapidinho, cheguei no caixa e tinha uns oito na fila. Calculei que não ia dar tempo, devolvi e fui embora sem comprar.

\textbf{Mateus:} O que acontece, na prática, entre você decidir vou levar isso e sair da loja com a peça?

\textbf{Felipe:} Eu meço mentalmente quanto tempo tenho, olho o tamanho da fila, e decido na hora se vale a pena esperar ou não.

\textbf{Mateus:} Isso acontece com frequência ou foi uma exceção?

\textbf{Felipe:} Acontece bastante, porque eu sempre vou com pouco tempo disponível.

\textbf{Mateus:} Você já usou algum tipo de autoatendimento, tipo mercado ou farmácia?

\textbf{Felipe:} Uso direto, prefiro até, é mais rápido que fila com atendente.

\textbf{Mateus:} Você confiaria em pagar por uma roupa pelo celular, sem falar com ninguém na loja?

\textbf{Felipe:} Confiaria numa boa, eu já uso pagamento por aproximação e pix pra tudo, não vejo problema nenhum.

\textbf{Mateus:} Já teve alguma experiência desconfortável com alarme ou checagem na saída?

\textbf{Felipe:} Nunca tive problema não.

\textbf{Mateus:} O que faria você confiar mais, ou menos, num processo assim?

\textbf{Felipe:} Confiaria mais se fosse rápido e sem burocracia. Menos se desse algum erro técnico que eu não soubesse resolver sozinho.

\textbf{Mateus:} Tem alguma coisa sobre comprar roupa que a gente não conversou e você acha importante?

\textbf{Felipe:} Acho que já cobrimos o principal pra mim, que é tempo. Tudo que eu reclamo tem a ver com perder tempo à toa.

\textbf{Mateus:} Se pudesse mudar uma coisa só nas lojas hoje, o que seria?

\textbf{Felipe:} Um caixa mais rápido, ou algum jeito de pagar sem precisar entrar na fila normal.

\textbf{Mateus:} Já viu em outro lugar alguma coisa parecida que resolvesse esse tipo de problema?

\textbf{Felipe:} Vi em loja de tecnologia lá fora que tem um totem de autopagamento, achei bem mais prático que fila comum.

\textbf{Mateus:} Você falou bastante sobre querer resolver o problema da fila. A Riachuelo está estudando justamente uma etiqueta de segurança na peça que se comunica com o aplicativo da loja. Ela sabe se a peça foi paga, e quando o pagamento é confirmado pelo aplicativo, ela se destrava sozinha, sem precisar passar pelo caixa. O que chama sua atenção nisso, pra bem ou pra mal?

\textbf{Felipe:} Achei muito bom, sinceramente, resolve exatamente o que mais me incomoda, que é perder tempo esperando pra pagar.

\textbf{Mateus:} Em que situação você usaria isso, e em que não usaria?

\textbf{Felipe:} Usaria sempre que possível, principalmente em compras rápidas tipo a que eu fiz semana passada.

\textbf{Mateus:} Isso resolve algum dos problemas que você mencionou antes, ou resolve outra coisa?

\textbf{Felipe:} Resolve exatamente o meu problema, que é fila. Não resolve a questão do app ser confuso, mas isso é outra parada.

\textbf{Mateus:} O que precisaria existir, ou o que teria que te garantir, para você confiar nesse processo?

\textbf{Felipe:} Precisava ser bem confiável tecnicamente, não travar, não dar erro de leitura. Se falhar uma vez já me deixa com pé atrás.

\textbf{Mateus:} Você entende isso como uma mudança só na trava da peça, ou como algo que muda a experiência de compra como um todo?

\textbf{Felipe:} Acho que muda bastante coisa, na verdade, porque se a peça já sabe que foi paga, dá pra imaginar até integrar com outras coisa, tipo avisar estoque em tempo real. Mas isso é só uma ideia minha.

\textbf{Mateus:} Isso mudaria a forma como você usa o vendedor hoje?

\textbf{Felipe:} Não muito, eu já uso pouco, então continuaria só recorrendo quando preciso de opinião sobre caimento.


\section{Entrevista 5}

\textbf{Danilo:} Como é sua relação com compra de roupa? Gosta, é necessidade, ou é um passeio?

\textbf{Nivaldo:} É mais por necessidade pontual. Eu não saio comprando roupa por lazer, mas também não é um sacrifício, eu até gosto de examinar o tecido, ver a qualidade, isso eu presto atenção.

\textbf{Danilo:} Com que frequência você compra, e onde costuma ir?

\textbf{Nivaldo:} Não é regular, depende da ocasião. Se tenho uma viagem ou um evento marcado, aí eu vou. Costumo ir na Zara do Iguatemi Campinas.

\textbf{Danilo:} Sozinho ou acompanhado?

\textbf{Nivaldo:} Geralmente sozinho. Já sou aposentado, tenho tempo, prefiro ir com calma no meu ritmo.

\textbf{Danilo:} Isso mudou nos últimos anos?

\textbf{Nivaldo:} Sim, desde que me aposentei tenho mais tempo, então quando vou, vou com mais calma do que quando trabalhava.

\textbf{Danilo:} Pensa na última vez que você comprou uma roupa numa loja física. Me conta essa história do início ao fim.

\textbf{Nivaldo:} Foi há uns dois mês, eu tinha uma viagem pro Nordeste marcada e precisava de roupa mais leve. Fui à loja, andei bastante pelos corredor, examinando os tecido, mexendo nas peças. Encontrei umas camisa de linho, experimentei, gostei do caimento de duas, e levei.

\textbf{Danilo:} Você já sabia o que estava procurando, ou foi mais explorando?

\textbf{Nivaldo:} Tinha uma ideia geral, camisa leve, mas fui bastante explorando até decidir qual exatamente.

\textbf{Danilo:} Como foi o momento de escolher e experimentar?

\textbf{Nivaldo:} Levei um tempo bom nisso, gosto de examinar bem antes de decidir, ver o tecido, sentir a textura.

\textbf{Danilo:} E o pagamento, como aconteceu?

\textbf{Nivaldo:} No cartão físico, no caixa. Não gosto muito de pagar pelo celular, prefiro o cartão na mão.

\textbf{Danilo:} Teve algum momento que você hesitou ou quase desistiu?

\textbf{Nivaldo:} Hesitei entre duas cor, mas não foi nada grave, decidi rápido.

\textbf{Danilo:} Alguma parte demorou mais do que você esperava?

\textbf{Nivaldo:} Não, achei o tempo todo dentro do esperado, porque eu já vou disposto a demorar um pouco.

\textbf{Danilo:} Isso é parecido com o normal pra você, ou foi diferente?

\textbf{Nivaldo:} É parecido. Eu costumo examinar bem antes de decidir, isso é característico de como eu compro.

\textbf{Danilo:} Quando você procura alguma coisa específica, como você faz pra achar?

\textbf{Nivaldo:} Ando pela loja com calma, olhando as seção. Não tenho pressa, então acabo encontrando por conta própria na maioria das vez.

\textbf{Danilo:} Já aconteceu de não encontrar o tamanho, a cor ou o modelo que queria?

\textbf{Nivaldo:} Já aconteceu algumas vez, principalmente tamanho, eu sou um pouco mais alto que a média.

\textbf{Danilo:} Você pesquisa em site ou aplicativo antes de ir à loja?

\textbf{Nivaldo:} Às vezes olho, mas confesso que não confio muito na precisão, prefiro conferir pessoalmente antes de decidir.

\textbf{Danilo:} O que você faz quando não acha o que quer, pede ajuda, procura mais, desiste?

\textbf{Nivaldo:} Procuro mais um pouco sozinho, e se realmente não encontrar, peço a um vendedor pra verificar no estoque.

\textbf{Danilo:} Isso te faz sair de mãos vazias com que frequência?

\textbf{Nivaldo:} Não é muito frequente, talvez uma vez a cada bom tempo.

\textbf{Danilo:} Se você soubesse, antes de sair de casa, que a loja tem o que você quer, isso mudaria alguma coisa?

\textbf{Nivaldo:} Ajudaria, sim, principalmente quando é uma peça específica que eu já sei que quero, evitaria uma ida à toa.

\textbf{Danilo:} Voltando à sua última compra, um vendedor te ajudou em algum momento?

\textbf{Nivaldo:} Ajudou no final, quando eu fui buscar um tamanho que não achei na araque.

\textbf{Danilo:} Tem alguma parte que você prefere resolver sozinho, sem falar com ninguém?

\textbf{Nivaldo:} A parte de escolher e examinar eu prefiro sozinho, gosto de decidir no meu tempo, sem alguém interferindo no meio.

\textbf{Danilo:} O que um vendedor resolve para você que você não resolveria sozinho?

\textbf{Nivaldo:} Buscar tamanho no estoque, principalmente, já que meu tamanho às vezes não está exposto.

\textbf{Danilo:} Já teve uma situação em que a ajuda de um vendedor atrapalhou mais do que ajudou, ou sentiu falta e não tinha ninguém?

\textbf{Nivaldo:} Já tive vendedor insistindo em me mostrar outras peça quando eu só queria olhar sozinho, isso me incomoda um pouco.

\textbf{Danilo:} Isso muda dependendo do tipo de peça ou do valor?

\textbf{Nivaldo:} Um pouco. Para peça mais cara eu até aceito mais uma segunda opinião, mas gosto de pedir eu mesmo, não que venham me abordar.

\textbf{Danilo:} Isso mudou com o tempo, você busca mais ou menos ajuda hoje?

\textbf{Nivaldo:} Acho que busco menos hoje, já tenho bastante experiência de comprar sozinho ao longo dos ano.

\textbf{Danilo:} Como foi o pagamento na sua última compra? Teve fila?

\textbf{Nivaldo:} Teve uma fila pequena, mas nada que me incomodasse muito.

\textbf{Danilo:} De modo geral, o que você acha do processo entre decidir levar a peça e sair da loja?

\textbf{Nivaldo:} Acho razoável. Não é algo que me tira do sério, desde que não seja excessivo.

\textbf{Danilo:} Já desistiu de comprar algo por causa da fila ou demora no caixa?

\textbf{Nivaldo:} Não me recordo de ter desistido por isso especificamente.

\textbf{Danilo:} O que acontece, na prática, entre você decidir vou levar isso e sair da loja com a peça?

\textbf{Nivaldo:} Eu levo a peça até o caixa, aguardo minha vez com tranquilidade, pago e saio.

\textbf{Danilo:} Isso acontece com frequência ou foi uma exceção?

\textbf{Nivaldo:} É basicamente sempre assim pra mim, não varia muito.

\textbf{Danilo:} Você já usou algum tipo de autoatendimento, tipo mercado ou farmácia?

\textbf{Nivaldo:} Já usei em supermercado algumas vez, funciona bem, embora eu ainda prefira o atendente em certas ocasião.

\textbf{Danilo:} Você confiaria em pagar por uma roupa pelo celular, sem falar com ninguém na loja?

\textbf{Nivaldo:} Tenho uma certa resistência, confesso. Prefiro o cartão físico, sinto mais controle sobre a transação.

\textbf{Danilo:} Já teve alguma experiência desconfortável com alarme ou checagem na saída?

\textbf{Nivaldo:} Não, nunca tive esse tipo de problema.

\textbf{Danilo:} O que faria você confiar mais, ou menos, num processo assim?

\textbf{Nivaldo:} Confiaria mais se houvesse um comprovante claro e imediato, e algum canal de contato caso algo desse errado.

\textbf{Danilo:} Tem alguma coisa sobre comprar roupa que a gente não conversou e você acha importante?

\textbf{Nivaldo:} Acho que a qualidade do tecido é algo que pouca gente valoriza hoje em dia, e pra mim isso pesa bastante na decisão.

\textbf{Danilo:} Se pudesse mudar uma coisa só nas lojas hoje, o que seria?

\textbf{Nivaldo:} Melhoraria a disponibilidade de tamanhos maiores, que é algo que sempre me limita.

\textbf{Danilo:} Já viu em outro lugar alguma coisa parecida que resolvesse esse tipo de problema?

\textbf{Nivaldo:} Não que eu me lembre agora.

\textbf{Danilo:} A Riachuelo está estudando uma etiqueta de segurança na peça que se comunica com o aplicativo da loja. Ela sabe se a peça foi paga, e quando o pagamento é confirmado pelo aplicativo, ela se destrava sozinha, sem precisar passar pelo caixa. O que chama sua atenção nisso, para o bem ou para o mal?

\textbf{Nivaldo:} Acho uma ideia interessante em termos de agilidade, mas tenho ressalvas quanto à segurança da transação e à minha própria resistência a pagar por celular.

\textbf{Danilo:} Em que situação você usaria isso, e em que situação não usaria?

\textbf{Nivaldo:} Talvez usasse para uma compra simples e de valor baixo. Para algo mais caro, ainda prefiro o processo tradicional, com um funcionário confirmando.

\textbf{Danilo:} Isso resolve algum dos problemas que você mencionou antes, ou resolve outra coisa?

\textbf{Nivaldo:} Não resolve minha questão principal, que é disponibilidade de tamanho. Resolve uma questão de agilidade no pagamento, que para mim não era prioridade.

\textbf{Danilo:} O que precisaria existir, ou o que teria que te garantir, para você confiar nesse processo?

\textbf{Nivaldo:} Um comprovante bem claro, e a certeza de que, em caso de erro, existe alguém disponível para resolver rapidamente.

\textbf{Danilo:} O que te preocuparia mais, a tecnologia falhar, ou perder contato com alguém que te ajuda quando precisa?

\textbf{Nivaldo:} As duas coisa me preocupam de forma parecida, mas se eu tivesse que escolher, diria que é a tecnologia falhar sem eu ter a quem recorrer.


\section{Entrevista 6}

\textbf{Lucas:} Bora, Thiago. Me conta, como é sua relação com compra de roupa? Curte comprar, é necessidade, ou rola mais como um passeio?

\textbf{Thiago:} Ah, eu curto sim, mas não fico enrolando não. Vou, olho rápido, se achar algo bom eu já levo.

\textbf{Lucas:} Com que frequência você compra, e onde costuma ir?

\textbf{Thiago:} Umas duas vez por mês eu dou uma olhada. Vou na Renner do Shopping Metrô Tatuapé, aqui perto de casa.

\textbf{Lucas:} Sozinho ou acompanhado, geralmente?

\textbf{Thiago:} Sozinho na maioria, às vezes com um amigo, mas cada um vê o seu, não fico esperando ele escolher.

\textbf{Lucas:} Isso mudou nos últimos anos?

\textbf{Thiago:} Mudou um pouco, hoje eu tenho menos tempo porque trabalho com aplicativo de carro e ainda faço faculdade, então quando eu vou é rapidinho mesmo.

\textbf{Lucas:} Boa. Pensa na última vez que você comprou uma roupa numa loja física. Conta essa história inteira, desde o que te fez ir até você sair.

\textbf{Thiago:} Foi semana passada. Precisava de uma camisa pra trabalhar, entrei, vi umas três opção, escolhi uma bem rápido, provei, servia, paguei e saí. Total, uns dez minuto.

\textbf{Lucas:} Você já sabia o que queria ou foi mais explorando?

\textbf{Thiago:} Sabia mais ou menos, mas acabei achando outra coisa que eu nem tinha pensado e gostei mais.

\textbf{Lucas:} Como foi a escolha e a experimentação?

\textbf{Thiago:} Rápido, testei ali mesmo, nem fui no provador, só na frente do espelho da araque mesmo.

\textbf{Lucas:} E o pagamento?

\textbf{Thiago:} No pix, pelo celular, sem fila, tranquilo.

\textbf{Lucas:} Teve algum momento que você hesitou ou quase desistiu?

\textbf{Thiago:} Não, nada disso, foi direto.

\textbf{Lucas:} Alguma parte demorou mais do que esperava?

\textbf{Thiago:} Não, foi tudo rápido mesmo.

\textbf{Lucas:} Isso é parecido com o normal pra você, ou foi diferente?

\textbf{Thiago:} É bem parecido, eu sempre vou rápido assim.

\textbf{Lucas:} Quando você procura alguma coisa específica, como você faz pra achar?

\textbf{Thiago:} Eu mesmo procuro, não gosto muito de ficar chamando vendedor toda hora.

\textbf{Lucas:} Já aconteceu de não achar o tamanho ou modelo que queria?

\textbf{Thiago:} Já, várias vez, principalmente tênis, que eu também compro lá às vezes.

\textbf{Lucas:} Você pesquisa em site ou aplicativo antes de ir?

\textbf{Thiago:} Às vezes vejo no Instagram da loja ou no site, mas geralmente vou direto mesmo ver pessoalmente.

\textbf{Lucas:} O que você faz quando não acha o que quer, pede ajuda, procura mais, desiste?

\textbf{Thiago:} Procuro um pouco, se não achar eu já vou embora, não fico esperando vendedor.

\textbf{Lucas:} Isso te faz sair de mãos vazias com que frequência?

\textbf{Thiago:} Direto, uma vez sim outra não.

\textbf{Lucas:} Se você soubesse antes que a loja tem o que você quer, isso mudaria alguma coisa?

\textbf{Thiago:} Mudaria total, eu ia direto sem perder tempo indo à toa.

\textbf{Lucas:} Voltando à sua última compra, um vendedor te ajudou em algum momento?

\textbf{Thiago:} Não, nem precisei, resolvi tudo sozinho.

\textbf{Lucas:} Tem alguma parte que você prefere resolver sozinho?

\textbf{Thiago:} Praticamente tudo. Não curto muito vendedor em cima ficando insistindo pra eu levar mais coisa.

\textbf{Lucas:} O que um vendedor resolve pra você que você não resolveria sozinho?

\textbf{Thiago:} Só quando realmente não acho um tamanho, aí eu pergunto se tem no estoque.

\textbf{Lucas:} Já teve uma situação em que a ajuda de um vendedor atrapalhou mais do que ajudou, ou sentiu falta e não tinha ninguém?

\textbf{Thiago:} Já tive vendedor muito em cima, ficando falando de promoção o tempo todo, isso me irrita, eu só quero olhar em paz.

\textbf{Lucas:} Isso muda dependendo do tipo de peça ou do valor?

\textbf{Thiago:} Não muito não, eu sou sempre do jeito que sou, gosto de resolver rápido e sozinho.

\textbf{Lucas:} Isso mudou com o tempo?

\textbf{Thiago:} Acho que não mudou muito, sempre fui assim, direto ao ponto.

\textbf{Lucas:} Como foi o pagamento na sua última compra? Teve fila?

\textbf{Thiago:} Não teve, mas quando tem eu já fico impaciente, confesso.

\textbf{Lucas:} De modo geral, o que você acha do processo entre decidir levar a peça e sair da loja?

\textbf{Thiago:} Acho que é a parte mais chata quando tem fila. Já escolhi, já decidi, só quero ir embora.

\textbf{Lucas:} Já desistiu de comprar algo por causa da fila ou demora no caixa?

\textbf{Thiago:} Já, com certeza. Se a fila tá grande eu simplesmente não compro, deixo pra lá.

\textbf{Lucas:} Me conta de uma vez que isso aconteceu.

\textbf{Thiago:} Foi numa Black Friday, a fila do caixa dava volta na loja, eu só olhei e virei as costa, nem tentei.

\textbf{Lucas:} O que acontece, na prática, entre você decidir vou levar isso e sair da loja com a peça?

\textbf{Thiago:} Eu olho o tamanho da fila antes até de ir pro caixa, se for grande eu já penso se vale a pena.

\textbf{Lucas:} Isso acontece com frequência ou foi uma exceção?

\textbf{Thiago:} Acontece bastante, principalmente em promoção e final de semana.

\textbf{Lucas:} Você já usou algum tipo de autoatendimento, tipo mercado ou farmácia?

\textbf{Thiago:} Uso direto, acho muito melhor que ficar esperando atendente.

\textbf{Lucas:} Você confiaria em pagar por uma roupa pelo celular, sem falar com ninguém na loja?

\textbf{Thiago:} Com certeza, eu já pago tudo assim no dia a dia, pix, cartão no celular, essas coisa.

\textbf{Lucas:} Já teve alguma experiência desconfortável com alarme ou checagem na saída?

\textbf{Thiago:} Uma vez o alarme tocou sem motivo, mas o segurança já nem ligou muito, deixou eu passar.

\textbf{Lucas:} O que faria você confiar mais, ou menos, num processo assim?

\textbf{Thiago:} Confiaria mais se fosse rápido, sem burocracia nenhuma.

\textbf{Lucas:} Tem alguma coisa sobre comprar roupa que a gente não conversou e você acha importante?

\textbf{Thiago:} Acho que já falei o principal, que é eu odiar perder tempo à toa numa loja.

\textbf{Lucas:} Se pudesse mudar uma coisa só nas lojas hoje, o que seria?

\textbf{Thiago:} Tiraria a fila do caixa, ou pelo menos deixaria mais rápida.

\textbf{Lucas:} Já viu em outro lugar alguma coisa parecida que resolvesse esse tipo de problema?

\textbf{Thiago:} Já vi loja de tênis que tem um totem de autopagamento, achei bem mais rápido que fila normal.

\textbf{Lucas:} Você já comprou pelo site alguma vez?

\textbf{Thiago:} Já, uso bastante inclusive, mas às vezes prefiro loja física quando quero levar na hora.

\textbf{Lucas:} Legal. Olha só, a Riachuelo está estudando uma etiqueta de segurança na peça que se comunica com o aplicativo da loja. Ela sabe se a peça foi paga, e quando o pagamento é confirmado pelo aplicativo, ela se destrava sozinha, sem precisar passar pelo caixa. O que você acha disso, pra bem ou pra mal?

\textbf{Thiago:} Achei sensacional, sinceramente. Resolve exatamente o que eu mais odeio, que é fila.

\textbf{Lucas:} Em que situação você usaria isso, e em que não usaria?

\textbf{Thiago:} Usaria sempre, não vejo motivo pra não usar, seria bem mais rápido pra mim.

\textbf{Lucas:} Isso resolve algum dos problemas que você mencionou antes, ou resolve outra coisa?

\textbf{Thiago:} Resolve o problema da fila com certeza, que é o que mais me incomoda.

\textbf{Lucas:} O que precisaria existir, ou o que teria que te garantir, para você confiar nesse processo?

\textbf{Thiago:} Só precisa funcionar direito, não travar, não dar erro. Se funcionar eu confio numa boa.

\textbf{Lucas:} Você imagina algum problema com uma peça saber que foi manuseada ou onde está dentro da loja?

\textbf{Thiago:} Não vejo problema não, pra mim é só tecnologia funcionando, não ligo muito pra esse tipo de coisa.


\section{Entrevista 7}

\textbf{Lucas:} E aí, Renata, vamos começar. Me conta, como é a sua relação com compra de roupa? Você gosta de comprar, é mais necessidade, ou é tipo um passeio?

\textbf{Renata:} Eu gosto bastante, viu. Pra mim é quase um momento de cuidar de mim mesma. Eu costumo chegar na loja já com uma referência na cabeça, às vezes uma foto que eu salvei, um look que eu vi em alguém, e vou atrás daquilo.

\textbf{Lucas:} Com que frequência você compra, e onde costuma ir?

\textbf{Renata:} Umas três, quatro semana eu dou uma passada. Vou na Zara do Iguatemi São Paulo, na Faria Lima.

\textbf{Lucas:} Sozinha ou acompanhada, geralmente?

\textbf{Renata:} Depende. Às vezes sozinha, às vezes com uma amiga que também gosta de moda, a gente troca ideia.

\textbf{Lucas:} Isso mudou nos últimos anos?

\textbf{Renata:} Um pouco. Hoje eu chego mais preparada, com referência de Instagram e Pinterest, antes eu ia mais no impulso.

\textbf{Lucas:} Legal. Pensa na última vez que você comprou uma roupa numa loja física. Me conta essa história do início ao fim.

\textbf{Renata:} Foi há umas duas semana. Eu tinha visto um look no Instagram de uma blogueira, um conjunto meio alfaiataria, e fui atrás de algo parecido. Cheguei na loja, expliquei pra vendedora o que eu queria, ela é uma consultora que já me atende há um tempo, ela foi trazendo opção, eu fui provando, ela dava opinião do caimento, e a gente foi ajustando até achar a combinação certa.

\textbf{Lucas:} Você já sabia exatamente o que estava procurando, ou foi mais explorando com a ajuda dela?

\textbf{Renata:} Uma mistura das duas coisa. Eu tinha a referência, mas ela foi essencial pra traduzir aquilo pro que a loja realmente tinha disponível.

\textbf{Lucas:} Como foi o momento de escolher e experimentar?

\textbf{Renata:} Foi bem prazeroso, na verdade. Passei bastante tempo no provador, ela ficava do lado de fora trazendo variação, e eu ia dando meu feedback.

\textbf{Lucas:} E o pagamento, como aconteceu?

\textbf{Renata:} No cartão, sem fila, porque essa loja é menor e mais tranquila que as de shopping popular.

\textbf{Lucas:} Teve algum momento que você hesitou ou quase desistiu?

\textbf{Renata:} Hesitei entre duas blazer, mas a vendedora ajudou a decidir, mostrando qual combinava melhor com o resto do look.

\textbf{Lucas:} Alguma parte demorou mais do que você esperava?

\textbf{Renata:} Demorou mais que o normal, mas porque eu quis, gostei de dedicar tempo pra isso.

\textbf{Lucas:} Isso é parecido com o que normalmente acontece pra você, ou foi diferente?

\textbf{Renata:} É bem parecido, eu sempre gosto de dedicar tempo e ter uma consultora ajudando.

\textbf{Lucas:} Quando você procura alguma coisa específica, como você faz pra achar?

\textbf{Renata:} Eu já vou direto falar com a vendedora que me conhece, explico o que eu quero, às vezes mostro a foto de referência no celular.

\textbf{Lucas:} Já aconteceu de não encontrar o tamanho, a cor ou o modelo que queria?

\textbf{Renata:} Já, mais de uma vez. Às vezes o modelo que eu vi na rede social nem chegou a essa loja específica.

\textbf{Lucas:} Você pesquisa em algum canal antes de ir à loja?

\textbf{Renata:} Sim, sempre. Redes sociais principalmente, às vezes o site pra ver se a peça existe antes de ir.

\textbf{Lucas:} O que você faz quando não acha o que quer, pede ajuda, procura mais, desiste?

\textbf{Renata:} Peço ajuda direto pra vendedora, ela costuma saber se tem em outra loja ou se pode pedir transferência.

\textbf{Lucas:} Isso te faz sair de mãos vazias com que frequência?

\textbf{Renata:} Acontece, talvez uma em cada quatro vez que eu vou atrás de algo específico.

\textbf{Lucas:} Se você soubesse, antes de sair de casa, que a loja tem o que você quer, isso mudaria alguma coisa?

\textbf{Renata:} Mudaria, com certeza. Eu economizaria idas frustradas, principalmente quando é algo que eu vi numa rede social e fiquei animada.

\textbf{Lucas:} Voltando à sua última compra, em algum momento a vendedora te ajudou de um jeito que você considera essencial?

\textbf{Renata:} Sim, o tempo todo, na verdade. Pra mim ela é praticamente uma consultora de estilo, não só alguém que me atende.

\textbf{Lucas:} Tem alguma parte da compra em que você prefere não falar com ninguém e resolver sozinha?

\textbf{Renata:} Olha, sinceramente, não muito. Eu gosto da companhia dela durante quase todo o processo. Talvez só na hora de decidir o pagamento em si, isso eu resolvo sozinha.

\textbf{Lucas:} O que uma vendedora resolve pra você que você não resolveria sozinha?

\textbf{Renata:} Ela tem uma visão de conjunto que eu não tenho, sabe combinar peça, sugerir algo que eu nem tinha pensado.

\textbf{Lucas:} Já teve uma situação em que a ajuda de uma vendedora atrapalhou mais do que ajudou, ou sentiu falta e não tinha ninguém?

\textbf{Renata:} Já tive vendedora que não tinha muito preparo, ficava só dizendo que tudo ficava bom, sem critério nenhum, isso me frustra porque eu quero uma opinião de verdade.

\textbf{Lucas:} Isso muda dependendo do tipo de peça ou do valor?

\textbf{Renata:} Muda um pouco, pra peça de valor mais alto eu quero ainda mais essa consultoria, pra básico eu resolvo mais rápido.

\textbf{Lucas:} Isso mudou com o tempo, você busca mais ou menos ajuda hoje do que buscava antes?

\textbf{Renata:} Acho que busco mais hoje, porque eu encontrei essa vendedora específica que entende meu gosto, isso fez eu confiar mais em pedir ajuda.

\textbf{Lucas:} Como foi o momento de pagar na sua última compra? Teve fila?

\textbf{Renata:} Não teve, essa loja é mais tranquila nesse sentido.

\textbf{Lucas:} De modo geral, o que você acha do processo entre decidir que vai levar uma peça e efetivamente sair da loja com ela?

\textbf{Renata:} Pra mim não é algo que pesa muito, porque raramente eu enfrento fila grande nessa loja que frequento.

\textbf{Lucas:} Já desistiu de comprar algo por causa da fila ou da demora no caixa?

\textbf{Renata:} Já, mas foi numa loja diferente, mais popular, onde eu fui uma vez num evento promocional e a fila estava enorme.

\textbf{Lucas:} O que acontece, na prática, entre você decidir vou levar isso e você sair da loja com a peça em mãos?

\textbf{Renata:} Normalmente é rápido pra mim, a vendedora já embrulha, eu pago, e saio tranquila.

\textbf{Lucas:} Você já usou algum tipo de autoatendimento, supermercado, farmácia, cinema, sem passar por um funcionário?

\textbf{Renata:} Uso sim, em cinema principalmente, acho prático.

\textbf{Lucas:} Você confiaria em pagar por uma peça de roupa pelo celular, sem precisar falar com ninguém na loja?

\textbf{Renata:} Confiaria em termos técnicos, eu uso bastante tecnologia no dia a dia. Mas confesso que sentiria falta da interação, não é só uma questão de confiança na segurança.

\textbf{Lucas:} Como você se sente ao sair de uma loja com uma compra, já teve alguma experiência desconfortável com alarme, checagem de nota ou verificação na saída?

\textbf{Renata:} Nunca tive problema, mas imagino que seria bem desconfortável, principalmente numa loja onde as pessoas te conhecem.

\textbf{Lucas:} O que faria você confiar mais, ou menos, num processo assim?

\textbf{Renata:} Confiaria mais se fosse algo opcional, que não substituísse a possibilidade de ter uma pessoa por perto quando eu quiser.

\textbf{Lucas:} Tem alguma coisa sobre comprar roupa em loja física que a gente não conversou e que você acha importante?

\textbf{Renata:} Acho que vale falar da questão de achar peça que eu vi online. Às vezes é frustrante ir atrás de algo específico e a loja física simplesmente não ter, mesmo que apareça na rede social da marca.

\textbf{Lucas:} Se você pudesse mudar uma única coisa em como as lojas funcionam hoje, o que seria?

\textbf{Renata:} Melhoraria essa integração entre o que aparece nas redes e o que realmente tem disponível fisicamente.

\textbf{Lucas:} Já viu, em outra loja, outro país ou outro tipo de serviço, alguma coisa que resolvesse um problema parecido?

\textbf{Renata:} Já vi loja lá fora que tinha uma tela mostrando o estoque em tempo real da unidade, achei bem interessante.

\textbf{Lucas:} Voltando no que você comentou sobre integrar o digital com o que tem na loja, a Riachuelo está estudando uma etiqueta de segurança na peça que passa a se comunicar com o aplicativo da loja. Ela sabe se a peça foi paga e, quando o pagamento é confirmado pelo aplicativo, ela se destrava sozinha, sem precisar passar pelo caixa. O que chama sua atenção nisso, para o bem ou para o mal?

\textbf{Renata:} Achei curioso. Pra bem, é interessante pensar que a peça passa a ter uma espécie de identidade, isso pode abrir portas pra outras coisa, tipo saber estoque em tempo real, que é o que eu comentei. Pra mal, fico pensando se isso não vai deixar a experiência mais fria, menos humana.

\textbf{Lucas:} Em que situação você usaria isso, e em que situação não usaria?

\textbf{Renata:} Usaria talvez pra uma reposição rápida, algo que eu já sei que quero. Não usaria numa compra como a que descrevi, que envolve toda uma consultoria, ali eu quero a vendedora do meu lado até o fim.

\textbf{Lucas:} Isso resolve algum dos problemas que você mencionou antes, ou resolve outra coisa?

\textbf{Renata:} Não resolve meu problema principal, que é achar a peça certa e ter uma boa consultoria. Resolve uma questão de agilidade que, sinceramente, não é minha maior dor.

\textbf{Lucas:} O que precisaria existir, ou o que teria que te garantir, para você confiar nesse processo?

\textbf{Renata:} Precisaria continuar tendo a opção de ter alguém por perto, mesmo usando essa tecnologia. Não queria que isso substituísse a vendedora, e sim que fosse um complemento.

\textbf{Lucas:} Você entende isso como uma mudança só na trava da peça, ou como algo que muda a experiência de compra como um todo?

\textbf{Renata:} Acho que pode mudar mais coisa, principalmente se der pra usar essa identidade da peça pra outras coisa, tipo saber se tem estoque, mas isso depende de como implementarem.

\textbf{Lucas:} Isso mudaria a forma como você usa a vendedora hoje?

\textbf{Renata:} Não, eu continuaria usando do mesmo jeito, isso resolveria uma parte técnica, mas não a parte que eu mais valorizo, que é a relação com ela.


\section{Entrevista 8}

\textbf{Danilo:} Como é sua relação com compra de roupa? Gosta de comprar, é necessidade, ou é um passeio?

\textbf{Camila:} É bem um passeio pra mim. Costumo ir com uma amiga no fim de semana, a gente fica horas andando, provando roupa, é mais um programa do que uma obrigação.

\textbf{Danilo:} Com que frequência você compra, e onde costuma ir?

\textbf{Camila:} Vou umas duas vezes por mês, mas nem sempre compro, às vezes é só pra passear mesmo. Vou na Renner do Shopping Villa-Lobos, aqui perto de casa.

\textbf{Danilo:} Sozinha ou acompanhada, geralmente?

\textbf{Camila:} Quase sempre com essa amiga. Sozinha eu vou mais raramente, só quando é algo urgente.

\textbf{Danilo:} Isso mudou nos últimos anos?

\textbf{Camila:} Não mudou muito, sempre gostei desse programa de ir passear em loja, isso é bem constante pra mim.

\textbf{Danilo:} Pensa na última vez que você comprou uma roupa numa loja física. Me conta essa história do início ao fim.

\textbf{Camila:} Foi um sábado, fui com minha amiga, a gente ficou tipo duas horas na loja. Eu não tinha nada específico em mente, fui vendo o que tinha, experimentando várias coisa só pra ver como ficava, sem compromisso de levar. No final acabei levando um vestido que eu nem tinha visto de cara, achei olhando com calma numa araque mais no fundo da loja.

\textbf{Danilo:} Você já sabia o que estava procurando, ou foi mais explorando?

\textbf{Camila:} Foi bem explorando, eu nem tinha ido com objetivo nenhum, só queria dar uma olhada.

\textbf{Danilo:} Como foi o momento de escolher e experimentar?

\textbf{Camila:} Provei umas seis, sete peça, foi bem demorado, mas eu gosto disso, faz parte da diversão.

\textbf{Danilo:} E o pagamento, como aconteceu?

\textbf{Camila:} No cartão, teve uma filinha pequena, mas nada que me incomodasse, a gente ainda ficou conversando na fila.

\textbf{Danilo:} Teve algum momento que você hesitou ou quase desistiu?

\textbf{Camila:} Hesitei entre o vestido e uma saia, mas no fim decidi pelo vestido.

\textbf{Danilo:} Alguma parte demorou mais do que você esperava?

\textbf{Camila:} Não, pra mim demorar faz parte, eu não vou com pressa então não incomoda.

\textbf{Danilo:} Isso é parecido com o que normalmente acontece pra você, ou foi diferente?

\textbf{Camila:} É bem parecido, eu sempre curto esse tempo de explorar a loja com calma.

\textbf{Danilo:} Quando você está procurando alguma coisa específica, como você faz para encontrar?

\textbf{Camila:} Eu ando por toda a loja, gosto de ver tudo, às vezes acho coisa legal em lugar que eu nem esperava.

\textbf{Danilo:} Já aconteceu de você não encontrar o tamanho, a cor ou o modelo que queria?

\textbf{Camila:} Já, às vezes o meu tamanho falta, principalmente em peça mais na moda.

\textbf{Danilo:} Você pesquisa em algum canal, site, aplicativo, redes sociais, antes de ir à loja?

\textbf{Camila:} Às vezes vejo alguma coisa no Instagram que me chama atenção, mas geralmente não é o motivo de eu ir, vou mais pela experiência mesmo.

\textbf{Danilo:} O que você faz quando isso acontece, pede ajuda, procura sozinha por mais tempo, desiste?

\textbf{Camila:} Procuro mais um pouco, e se não achar eu deixo pra lá, não é o fim do mundo, sempre tem outra coisa que eu gosto.

\textbf{Danilo:} Isso te faz sair da loja de mãos vazias com que frequência?

\textbf{Camila:} Às vezes eu saio sem comprar, mas não é porque não achei, às vezes é só porque não tinha nada que eu quisesse levar mesmo.

\textbf{Danilo:} Se você soubesse, antes de sair de casa, que a loja tem o que você quer, isso mudaria alguma coisa em como você compra?

\textbf{Camila:} Um pouco, mas pra mim a graça é justamente não saber, ir descobrir. Se eu já soubesse tudo antes, perderia a parte de surpresa.

\textbf{Danilo:} Voltando à sua última compra, em algum momento um vendedor te ajudou? Como foi isso?

\textbf{Camila:} Só no final, quando eu fui pagar, uma moça perguntou se eu queria ajuda com o tamanho de uma peça que eu tinha deixado de lado, mas eu já tinha decidido sozinha.

\textbf{Danilo:} Tem alguma parte da compra em que você prefere não falar com ninguém e resolver sozinha?

\textbf{Camila:} A parte de escolher e experimentar, com certeza. Gosto de fazer isso no meu ritmo, sem ninguém em cima.

\textbf{Danilo:} O que um vendedor resolve para você que você não resolveria sozinha?

\textbf{Camila:} Achar tamanho que sumiu da araque, isso sim eu peço ajuda.

\textbf{Danilo:} Já teve uma situação em que a ajuda de um vendedor atrapalhou mais do que ajudou, ou o contrário?

\textbf{Camila:} Já tive vendedor meio em cima demais, perguntando toda hora se eu precisava de algo, isso atrapalha o clima de passear com calma.

\textbf{Danilo:} Isso muda dependendo do tipo de peça ou do valor?

\textbf{Camila:} Não muito, eu sou sempre parecida, gosto de explorar sozinha independente do valor da peça.

\textbf{Danilo:} Isso mudou com o tempo, você busca mais ou menos ajuda hoje do que buscava antes?

\textbf{Camila:} Acho que sempre fui assim, não mudou muito não.

\textbf{Danilo:} Como foi o momento de pagar na sua última compra? Teve fila?

\textbf{Camila:} Teve uma filinha, mas rápida, uns cinco minuto.

\textbf{Danilo:} De modo geral, o que você acha do processo entre decidir que vai levar uma peça e efetivamente sair da loja com ela?

\textbf{Camila:} Normalmente não me incomoda, só fica chato mesmo em datas especiais, tipo Black Friday, aí sim vira um problema.

\textbf{Danilo:} Já desistiu de comprar algo por causa da fila ou da demora no caixa?

\textbf{Camila:} Já, uma vez na Black Friday, a fila tava gigante, desisti e fui embora sem levar nada.

\textbf{Danilo:} O que acontece, na prática, entre você decidir vou levar isso e você sair da loja com a peça em mãos?

\textbf{Camila:} Normalmente é tranquilo, eu vou até o caixa, espero um pouco, pago e saio conversando com minha amiga.

\textbf{Danilo:} Isso acontece com frequência ou foi uma exceção?

\textbf{Camila:} A fila grande é mais exceção, coisa de data especial mesmo.

\textbf{Danilo:} Você já usou algum tipo de autoatendimento, supermercado, farmácia, cinema, sem passar por um funcionário?

\textbf{Camila:} Já usei em cinema e em algumas loja de roupa que já têm isso, achei bem prático quando funciona direito.

\textbf{Danilo:} Você confiaria em pagar por uma peça de roupa pelo celular, sem precisar falar com ninguém na loja?

\textbf{Camila:} Confiaria, eu já uso bastante esse tipo de coisa no dia a dia, não vejo grande problema.

\textbf{Danilo:} Como você se sente ao sair de uma loja com uma compra, já teve alguma experiência desconfortável com alarme, checagem de nota ou verificação na saída?

\textbf{Camila:} Nunca tive problema, mas imagino que seria bem chato passar por isso.

\textbf{Danilo:} O que faria você confiar mais, ou menos, num processo assim?

\textbf{Camila:} Confiaria mais se fosse simples de usar e não desse trabalho nenhum extra.

\textbf{Danilo:} Tem alguma coisa sobre comprar roupa em loja física que a gente não conversou e que você acha importante?

\textbf{Camila:} Acho que a ambientação da loja importa bastante pra mim, música, iluminação, isso faz parte da experiência de passear.

\textbf{Danilo:} Se você pudesse mudar uma única coisa em como as lojas funcionam hoje, o que seria?

\textbf{Camila:} Deixaria os provador mais confortável, com espelho melhor e mais espaço.

\textbf{Danilo:} Já viu, em outra loja, outro país ou outro tipo de serviço, alguma coisa que resolvesse um problema parecido com os que você mencionou?

\textbf{Camila:} Já vi loja com provador que tira foto e mostra numa tela como ficou de frente e de costas, achei bem legal.

\textbf{Danilo:} A Riachuelo está estudando uma etiqueta de segurança na peça que se comunica com o aplicativo da loja. Ela sabe se a peça foi paga e, quando o pagamento é confirmado pelo aplicativo, ela se destrava sozinha, sem precisar passar pelo caixa. O que chama sua atenção nisso, para o bem ou para o mal?

\textbf{Camila:} Achei legal a ideia, principalmente pra quando eu já sei o que quero e não preciso ficar na fila. Mas confesso que não é meu maior problema hoje.

\textbf{Danilo:} Em que situação você usaria isso, e em que situação não usaria?

\textbf{Camila:} Usaria numa compra rápida, tipo reposição de básico. Não usaria quando eu tô no clima de passear, ali eu nem me importo com a fila.

\textbf{Danilo:} Isso resolve algum dos problemas que você mencionou antes, ou resolve outra coisa?

\textbf{Camila:} Resolve a questão da fila em data especial, que é a única vez que isso realmente me incomoda.

\textbf{Danilo:} O que precisaria existir, ou o que teria que te garantir, para você confiar nesse processo?

\textbf{Camila:} Precisava ser simples de usar, sem burocracia, e ter algum aviso claro confirmando que deu tudo certo.

\textbf{Danilo:} O que te preocuparia mais, a tecnologia falhar, ou perder contato com alguém que te ajuda quando precisa?

\textbf{Camila:} Acho que nenhuma das duas me preocupa muito, porque eu já uso pouco vendedor mesmo. Talvez a tecnologia falhar seja um pouco mais chato.


\section{Entrevista 9}

\textbf{Mateus:} Como é a sua relação com compra de roupa? Gosta de comprar, é necessidade, ou funciona como um passeio?

\textbf{Débora:} Um pouco das duas coisa. Eu tenho um salão de beleza, então preciso estar sempre bem apresentada, é meio profissional isso pra mim, mas também gosto do momento de comprar, não é só obrigação.

\textbf{Mateus:} Com que frequência você compra, e onde costuma ir?

\textbf{Débora:} Vou umas três semana, sempre na Renner do Shopping Metrô Santa Cruz, na Vila Mariana, perto do meu salão.

\textbf{Mateus:} Sozinha ou acompanhada, geralmente?

\textbf{Débora:} Sozinha na maioria das vezes, meu tempo é curto, não dá pra combinar com ninguém.

\textbf{Mateus:} Isso mudou nos últimos anos?

\textbf{Débora:} Mudou, desde que o salão cresceu eu tenho menos tempo, então preciso ser mais objetiva quando vou comprar.

\textbf{Mateus:} Pensa na última vez que você comprou uma roupa numa loja física. Me conta essa história do início ao fim.

\textbf{Débora:} Foi há uns dez dia. Eu precisava de uma roupa pra um evento do setor, fui na loja no meu horário de almoço. Já fui direto falar com a vendedora que sempre me atende, ela já sabe meu gosto, trouxe duas opção, eu experimentei rápido e decidi por uma.

\textbf{Mateus:} Me conta mais sobre essa vendedora que já te conhece, como foi que isso começou?

\textbf{Débora:} Ah, é a Cristiane, ela trabalha lá tem uns dois ano. Depois de umas três vez que ela me atendeu, ela começou a entender meu estilo, o que funciona pro meu corpo, e isso facilitou muito, hoje eu confio bastante no que ela sugere.

\textbf{Mateus:} Nesse dia, você já sabia o que estava procurando, ou foi mais ela que guiou?

\textbf{Débora:} Foi uma mistura. Eu disse que precisava de algo pra evento formal, ela que trouxe as opção específica.

\textbf{Mateus:} Como foi o momento de escolher e experimentar?

\textbf{Débora:} Rápido, porque ela já trouxe coisa que combinava comigo, eu não precisei ficar testando muita coisa.

\textbf{Mateus:} E o pagamento, como aconteceu?

\textbf{Débora:} No cartão, sem fila, porque eu fui num horário mais vazio de propósito.

\textbf{Mateus:} Teve algum momento que você hesitou ou quase desistiu?

\textbf{Débora:} Não, dessa vez foi bem tranquilo.

\textbf{Mateus:} Alguma parte demorou mais do que você esperava?

\textbf{Débora:} Não, foi dentro do que eu esperava, uns vinte minuto.

\textbf{Mateus:} Isso é parecido com o que normalmente acontece pra você, ou foi diferente?

\textbf{Débora:} É bem parecido, eu sempre tento ir em horário vazio e já busco a Cristiane direto.

\textbf{Mateus:} Quando você está procurando alguma coisa específica, como você faz para encontrar?

\textbf{Débora:} Eu já vou direto falar com ela, não fico procurando sozinha, prefiro economizar tempo assim.

\textbf{Mateus:} Já aconteceu de você não encontrar o tamanho, a cor ou o modelo que queria?

\textbf{Débora:} Já, algumas vez, principalmente quando é uma peça mais nova, que ainda não chegou em quantidade.

\textbf{Mateus:} Você pesquisa em algum canal, site, aplicativo, redes sociais, antes de ir à loja?

\textbf{Débora:} Quase nunca, prefiro deixar pra Cristiane resolver, ela sabe melhor que eu o que tem disponível.

\textbf{Mateus:} O que você faz quando isso acontece, pede ajuda, procura sozinha por mais tempo, desiste?

\textbf{Débora:} Peço pra ela verificar em outra loja ou pedir transferência, quase nunca eu desisto de vez.

\textbf{Mateus:} Isso te faz sair da loja de mãos vazias com que frequência?

\textbf{Débora:} Raramente, porque ela geralmente resolve de algum jeito.

\textbf{Mateus:} Se você soubesse, antes de sair de casa, que a loja tem o que você quer, isso mudaria alguma coisa em como você compra?

\textbf{Débora:} Mudaria um pouco o planejamento, mas hoje eu já confio tanto na Cristiane que isso pesa menos pra mim.

\textbf{Mateus:} Voltando à sua última compra, em algum momento um vendedor te ajudou? Como foi isso?

\textbf{Débora:} O tempo todo, como eu já falei, pra mim ela é essencial no processo.

\textbf{Mateus:} Tem alguma parte da compra em que você prefere não falar com ninguém e resolver sozinha?

\textbf{Débora:} Só a decisão final de pagar mesmo, o resto eu gosto de fazer com ela do lado.

\textbf{Mateus:} O que um vendedor resolve para você que você não resolveria sozinha?

\textbf{Débora:} Ela filtra as opção pra mim, eu tenho pouco tempo, então isso é essencial, evita eu perder tempo vendo coisa que não combina.

\textbf{Mateus:} Já teve uma situação em que a ajuda de um vendedor atrapalhou mais do que ajudou, ou o contrário?

\textbf{Débora:} Já senti falta uma vez que ela estava de férias e outra vendedora me atendeu, não foi a mesma coisa, ela não conhecia meu gosto e sugeriu coisa que não combinava.

\textbf{Mateus:} Interessante, me conta mais sobre isso, como foi essa experiência?

\textbf{Débora:} Foi frustrante, sinceramente. Perdi mais tempo do que o normal porque eu tive que explicar tudo de novo pra alguém que não me conhecia, e ainda assim as sugestão não foram tão boas.

\textbf{Mateus:} Isso muda dependendo do tipo de peça ou do valor?

\textbf{Débora:} Um pouco, pra peça mais cara eu quero ainda mais certeza, então dependo mais dela nesses caso.

\textbf{Mateus:} Isso mudou com o tempo, você busca mais ou menos ajuda hoje do que buscava antes?

\textbf{Débora:} Busco mais hoje, porque eu encontrei alguém de confiança, antes eu tentava resolver mais sozinha e errava bastante.

\textbf{Mateus:} Como foi o momento de pagar na sua última compra? Teve fila?

\textbf{Débora:} Não teve, mas geralmente tem uma filinha no horário de almoço.

\textbf{Mateus:} De modo geral, o que você acha do processo entre decidir que vai levar uma peça e efetivamente sair da loja com ela?

\textbf{Débora:} Pra mim o que mais pesa é o tempo total gasto, incluindo achar a peça certa, não só o pagamento em si.

\textbf{Mateus:} Já desistiu de comprar algo por causa da fila ou da demora no caixa?

\textbf{Débora:} Já, uma vez o tempo de almoço tava acabando e eu tinha que voltar pro salão, deixei a peça e saí.

\textbf{Mateus:} O que acontece, na prática, entre você decidir vou levar isso e você sair da loja com a peça em mãos?

\textbf{Débora:} Normalmente é rápido, porque como eu falei, eu já vou direto com a Cristiane e ela agiliza tudo.

\textbf{Mateus:} Isso acontece com frequência ou foi uma exceção?

\textbf{Débora:} A fila grande é mais exceção, o normal é ser rápido pra mim.

\textbf{Mateus:} Você já usou algum tipo de autoatendimento, supermercado, farmácia, cinema, sem passar por um funcionário?

\textbf{Débora:} Já usei em farmácia, funciona bem pra coisa simples.

\textbf{Mateus:} Você confiaria em pagar por uma peça de roupa pelo celular, sem precisar falar com ninguém na loja?

\textbf{Débora:} Confiaria na parte técnica, mas sinceramente prefiro que a Cristiane esteja envolvida, confio mais na palavra dela do que num aplicativo.

\textbf{Mateus:} Como você se sente ao sair de uma loja com uma compra, já teve alguma experiência desconfortável com alarme, checagem de nota ou verificação na saída?

\textbf{Débora:} Nunca tive problema, mas confesso que ficaria incomodada se acontecesse, principalmente numa loja onde as pessoas me conhecem.

\textbf{Mateus:} O que faria você confiar mais, ou menos, num processo assim?

\textbf{Débora:} Confiaria mais se ainda envolvesse a Cristiane de alguma forma, tipo ela confirmando que deu tudo certo.

\textbf{Mateus:} Tem alguma coisa sobre comprar roupa em loja física que a gente não conversou e que você acha importante?

\textbf{Débora:} Acho que a relação de confiança com quem atende é a coisa mais importante pra mim, mais do que qualquer outra coisa que a gente conversou.

\textbf{Mateus:} Se você pudesse mudar uma única coisa em como as lojas funcionam hoje, o que seria?

\textbf{Débora:} Garantiria que meu horário de almoço sempre coincidisse com a Cristiane estar na loja, isso resolveria praticamente tudo pra mim.

\textbf{Mateus:} Já viu, em outra loja, outro país ou outro tipo de serviço, alguma coisa que resolvesse um problema parecido?

\textbf{Débora:} Não muito, não. Minha solução seria mais simples, tipo saber a escala dela com antecedência.

\textbf{Mateus:} Você comentou bastante sobre a importância da Cristiane no processo. A Riachuelo está estudando uma etiqueta de segurança na peça que se comunica com o aplicativo da loja. Ela sabe se a peça foi paga, e quando o pagamento é confirmado pelo aplicativo, ela se destrava sozinha, sem precisar passar pelo caixa. O que chama sua atenção nisso, para o bem ou para o mal?

\textbf{Débora:} Fico com um pé atrás, viu, porque isso parece tirar justamente a parte que eu mais valorizo, que é a pessoa envolvida no processo.

\textbf{Mateus:} Em que situação você usaria isso, e em que situação não usaria?

\textbf{Débora:} Talvez usasse pra uma reposição bem simples, tipo uma peça básica que eu já sei que quero. Pra roupa de evento, como a que comprei, eu ia querer a Cristiane do mesmo jeito.

\textbf{Mateus:} Isso resolve algum dos problemas que você mencionou antes, ou resolve outra coisa?

\textbf{Débora:} Não resolve meu problema principal, que é ter alguém de confiança disponível. Resolve uma questão de agilidade que não é minha prioridade.

\textbf{Mateus:} O que precisaria existir, ou o que teria que te garantir, para você confiar nesse processo?

\textbf{Débora:} Precisaria de alguma forma manter a vendedora envolvida, nem que fosse ela confirmando pelo sistema dela que está tudo certo.

\textbf{Mateus:} Você entende isso como uma mudança só na trava da peça, ou como algo que muda a experiência de compra como um todo?

\textbf{Débora:} Acho que, do jeito que você descreveu, parece mudar mais a parte final, o pagamento e a saída. O resto, escolher com a Cristiane, continuaria igual pra mim.


\section{Entrevista 10}

\textbf{Lucas:} E aí, Wagner, vamos começar. Me conta, como é sua relação com compra de roupa? Gosta de comprar, é necessidade, ou é tipo um passeio?

\textbf{Wagner:} É necessidade, com certeza. Eu quase nunca compro pra mim, é sempre pros meus filho, então já vou com aquela missão na cabeça.

\textbf{Lucas:} Com que frequência você compra, e onde costuma ir?

\textbf{Wagner:} Umas duas vezes por mês eu acabo indo, tenho três filho, sempre alguém precisa de alguma coisa. Vou na C\&A do Shopping Center Norte, aqui da Zona Norte.

\textbf{Lucas:} Sozinho ou acompanhado, geralmente?

\textbf{Wagner:} Sozinho na maioria das vezes, porque eu vou nos meus horário de folga, e nem sempre dá pra levar as criança.

\textbf{Lucas:} Isso mudou nos últimos anos?

\textbf{Wagner:} Mudou bastante, desde que nasceu o mais novo ficou mais corrido, e comprar roupa de criança é sempre correndo atrás do tamanho que eles cresce rápido.

\textbf{Lucas:} Pensa na última vez que você comprou uma roupa numa loja física. Me conta essa história do início ao fim.

\textbf{Wagner:} Foi semana passada, precisava comprar uma calça pro meu filho do meio, que tá numa fase de crescimento, acho que uns dois número por ano. Fui sozinho, sem ele, então tive que tentar adivinhar o tamanho, olhei uma calça que ele tinha usado recentemente pra comparar antes de sair de casa. Cheguei lá, procurei o tamanho, comparei com outras peça pra ver se batia, levei duas opção pra garantir.

\textbf{Lucas:} Você já sabia exatamente o que estava procurando, ou foi mais tentando adivinhar mesmo?

\textbf{Wagner:} Foi bem tentando adivinhar, sinceramente. Eu tirei uma foto da etiqueta de uma calça dele em casa antes de sair, mas mesmo assim eu fico inseguro porque cada marca tem um tamanho diferente.

\textbf{Lucas:} Como foi o momento de escolher?

\textbf{Wagner:} Comparei bastante, fiquei um tempo ali medindo com as mão mesmo, tentando ver se batia com o que eu lembrava dele.

\textbf{Lucas:} E o pagamento, como aconteceu?

\textbf{Wagner:} No cartão, teve uma filazinha, mas nada muito longo.

\textbf{Lucas:} Teve algum momento que você hesitou ou quase desistiu?

\textbf{Wagner:} Hesitei bastante entre dois tamanho, cheguei a levar os dois pra depois trocar se precisasse.

\textbf{Lucas:} Alguma parte demorou mais do que você esperava?

\textbf{Wagner:} Sim, a parte de decidir o tamanho, isso sempre demora mais pra mim do que eu queria.

\textbf{Lucas:} Isso é parecido com o que normalmente acontece pra você, ou foi diferente?

\textbf{Wagner:} É bem parecido, toda vez que eu vou comprar pros meus filho sem eles é essa insegurança do tamanho.

\textbf{Lucas:} Quando você está procurando alguma coisa específica, como você faz para encontrar?

\textbf{Wagner:} Eu ando procurando sozinho primeiro, comparando os número, e se eu fico muito em dúvida chamo alguém da loja pra confirmar.

\textbf{Lucas:} Já aconteceu de você não encontrar o tamanho, a cor ou o modelo que queria?

\textbf{Wagner:} Já, bastante, principalmente tamanho de criança que varia muito entre uma coleção e outra.

\textbf{Lucas:} Você pesquisa em algum canal, site, aplicativo, redes sociais, antes de ir à loja?

\textbf{Wagner:} Às vezes eu olho no site pra ver se tem o tamanho, mas confesso que eu não confio cem por cento, prefiro conferir na loja mesmo.

\textbf{Lucas:} O que você faz quando isso acontece, pede ajuda, procura sozinho por mais tempo, desiste?

\textbf{Wagner:} Peço ajuda pra vendedora, geralmente elas têm mais prática com tamanho de criança do que eu.

\textbf{Lucas:} Isso te faz sair da loja de mãos vazias com que frequência?

\textbf{Wagner:} De vez em quando, principalmente quando eu não confio em nenhum dos tamanho disponível e prefiro levar meu filho depois pra ele mesmo provar.

\textbf{Lucas:} Se você soubesse, antes de sair de casa, que a loja tem o que você quer, isso mudaria alguma coisa em como você compra?

\textbf{Wagner:} Mudaria bastante, eu ia com mais confiança, sem ficar com medo de comprar errado.

\textbf{Lucas:} Voltando à sua última compra, em algum momento um vendedor te ajudou? Como foi isso?

\textbf{Wagner:} Ajudou, sim, uma moça me deu uma dica de comparar com o comprimento do braço dele pra estimar melhor o tamanho, achei útil.

\textbf{Lucas:} Tem alguma parte da compra em que você prefere não falar com ninguém e resolver sozinho?

\textbf{Wagner:} A parte de escolher a cor e o modelo eu resolvo sozinho, sei bem o gosto dele. É só tamanho mesmo que eu preciso de ajuda.

\textbf{Lucas:} O que um vendedor resolve para você que você não resolveria sozinho?

\textbf{Wagner:} A questão do tamanho, com certeza, é onde eu mais erro sozinho.

\textbf{Lucas:} Já teve uma situação em que a ajuda de um vendedor atrapalhou mais do que ajudou, ou o contrário?

\textbf{Wagner:} Já precisei de ajuda e não tinha ninguém disponível, fiquei bom tempo esperando, isso é chato quando eu tô com pouco tempo de folga.

\textbf{Lucas:} Isso muda dependendo do tipo de peça ou do valor?

\textbf{Wagner:} Muda um pouco, pra tênis, por exemplo, eu sempre peço ajuda porque eu erro muito o número.

\textbf{Lucas:} Isso mudou com o tempo, você busca mais ou menos ajuda hoje do que buscava antes?

\textbf{Wagner:} Busco mais hoje, já aprendi que sozinho eu erro bastante e depois tenho que voltar pra trocar, o que dá mais trabalho ainda.

\textbf{Lucas:} Como foi o pagamento na sua última compra? Teve fila?

\textbf{Wagner:} Teve uma filinha pequena, nada muito demorado.

\textbf{Lucas:} De modo geral, o que você acha do processo entre decidir levar a peça e sair da loja?

\textbf{Wagner:} Pra mim o que mais pesa não é bem essa parte, é antes, é a insegurança de saber se o tamanho tá certo.

\textbf{Lucas:} Já desistiu de comprar algo por causa da fila ou demora no caixa?

\textbf{Wagner:} Não me lembro de ter desistido por isso especificamente, meu problema é mais outro.

\textbf{Lucas:} O que acontece, na prática, entre você decidir vou levar isso e sair da loja com a peça?

\textbf{Wagner:} Eu levo as duas opção que eu separei, vou até o caixa, pago as duas, com a intenção de trocar depois se precisar.

\textbf{Lucas:} Isso acontece com frequência ou foi uma exceção?

\textbf{Wagner:} Acontece direto, é praticamente sempre assim quando eu compro sem meus filho.

\textbf{Lucas:} Você já usou algum tipo de autoatendimento, supermercado, farmácia, cinema, sem passar por um funcionário?

\textbf{Wagner:} Já usei no mercado, funciona bem pra coisa simples que eu já sei o que é.

\textbf{Lucas:} Você confiaria em pagar por uma peça de roupa pelo celular, sem precisar falar com ninguém na loja?

\textbf{Wagner:} Confiaria na parte de pagar, isso não me preocupa. O que me preocupa mesmo é sempre a questão do tamanho, isso nenhum aplicativo resolve pra mim.

\textbf{Lucas:} Como você se sente ao sair de uma loja com uma compra, já teve alguma experiência desconfortável com alarme, checagem de nota ou verificação na saída?

\textbf{Wagner:} Nunca tive problema não.

\textbf{Lucas:} O que faria você confiar mais, ou menos, num processo assim?

\textbf{Wagner:} Confiaria mais se fosse simples e rápido, eu não vejo muito problema nessa parte.

\textbf{Lucas:} Tem alguma coisa sobre comprar roupa em loja física que a gente não conversou e que você acha importante?

\textbf{Wagner:} Acho que devia ter algum jeito melhor de saber o tamanho certo sem precisar levar a criança junto toda vez.

\textbf{Lucas:} Se você pudesse mudar uma única coisa em como as lojas funcionam hoje, o que seria?

\textbf{Wagner:} Um jeito de comparar tamanho de forma mais fácil, tipo alguma tabela melhor ou algo que me ajudasse a acertar de primeira.

\textbf{Lucas:} Já viu, em outra loja, outro país ou outro tipo de serviço, alguma coisa que resolvesse um problema parecido com os que você mencionou?

\textbf{Wagner:} Não, não vi nada assim não.

\textbf{Lucas:} A Riachuelo está estudando uma etiqueta de segurança na peça que se comunica com o aplicativo da loja. Ela sabe se a peça foi paga, e quando o pagamento é confirmado pelo aplicativo, ela se destrava sozinha, sem precisar passar pelo caixa. O que chama sua atenção nisso, para o bem ou para o mal?

\textbf{Wagner:} Achei ok, mas sinceramente não é isso que resolveria meu maior problema, que é acertar o tamanho.

\textbf{Lucas:} Em que situação você usaria isso, e em que situação não usaria?

\textbf{Wagner:} Usaria pra qualquer compra, eu não vejo desvantagem, mas também não vejo como uma solução pro que mais me incomoda.

\textbf{Lucas:} Isso resolve algum dos problemas que você mencionou antes, ou resolve outra coisa?

\textbf{Wagner:} Resolve uma coisa que nem era meu problema principal. Meu problema é tamanho, isso aqui resolve fila e pagamento.

\textbf{Lucas:} O que precisaria existir, ou o que teria que te garantir, para você confiar nesse processo?

\textbf{Wagner:} Só precisaria funcionar direito, sem complicação, aí eu uso numa boa.

\textbf{Lucas:} O que te preocuparia mais, a tecnologia falhar, ou perder contato com alguém que te ajuda quando precisa?

\textbf{Wagner:} Perder contato com quem ajuda, porque é dessa ajuda que eu mais preciso, mais do que da tecnologia em si.


\section{Entrevista 11}

\textbf{Danilo:} Como é sua relação com compra de roupa? Gosta de comprar, é necessidade, ou é um passeio?

\textbf{Luana:} É complicado responder isso de forma simples. Eu estudo moda, então entrar numa loja pra mim é quase um trabalho de observação, eu vou mais pra ver o que estão fazendo, tendência, do que necessariamente pra comprar.

\textbf{Danilo:} Com que frequência você compra, e onde costuma ir?

\textbf{Luana:} Comprar mesmo, pouquíssimas vezes, talvez a cada dois, três mês. Mas entrar numa loja pra olhar, isso eu faço toda semana, várias loja diferente, tipo a Zara do Shopping Pátio Higienópolis, mas não só ela.

\textbf{Danilo:} Sozinha ou acompanhada, geralmente?

\textbf{Luana:} Sozinha, sempre. Gosto de prestar atenção sem ninguém puxando papo do lado.

\textbf{Danilo:} Isso mudou nos últimos anos?

\textbf{Luana:} Mudou desde que eu entrei na faculdade. Antes eu comprava mais por impulso, hoje eu observo muito mais do que compro.

\textbf{Danilo:} Pensa na última vez que você comprou uma roupa numa loja física. Me conta essa história do início ao fim.

\textbf{Luana:} Foi há uns três mês, acho. Eu tava andando pela loja sem intenção nenhuma de comprar, só olhando as peça nova da coleção, e vi uma jaqueta com um corte diferente que me chamou atenção. Fui no provador mais por curiosidade de ver como ficava do que já pensando em levar, gostei, e levei.

\textbf{Danilo:} Você já sabia o que estava procurando, ou foi mais explorando?

\textbf{Luana:} Totalmente explorando, eu nem tinha ido com intenção de comprar aquele dia.

\textbf{Danilo:} Como foi o momento de escolher e experimentar?

\textbf{Luana:} Rápido, na real, porque quando eu já sei que gostei do corte, eu não fico enrolando muito.

\textbf{Danilo:} E o pagamento, como aconteceu?

\textbf{Luana:} No pix, sem fila, porque eu fui num horário bem vazio de manhã.

\textbf{Danilo:} Teve algum momento que você hesitou ou quase desistiu?

\textbf{Luana:} Hesitei um pouco pelo preço, porque geralmente eu prefiro comprar em brechó, que é bem mais em conta e mais sustentável, mas como era uma peça diferente, decidi levar mesmo assim.

\textbf{Danilo:} Alguma parte demorou mais do que você esperava?

\textbf{Luana:} Não, foi tudo rápido.

\textbf{Danilo:} Isso é parecido com o que normalmente acontece pra você, ou foi diferente?

\textbf{Luana:} Foi diferente, porque normalmente eu nem compro, só observo. Comprar de fato é raro.

\textbf{Danilo:} Quando você está procurando alguma coisa específica, como você faz para encontrar?

\textbf{Luana:} Raramente eu vou atrás de algo específico, mas quando vou, eu prefiro procurar sozinha, gosto de tocar o tecido, ver a qualidade antes de qualquer coisa.

\textbf{Danilo:} Já aconteceu de você não encontrar o tamanho, a cor ou o modelo que queria?

\textbf{Luana:} Já, mas como eu compro pouco, isso não me incomoda tanto quanto eu acho que incomodaria outra pessoa.

\textbf{Danilo:} Você pesquisa em algum canal, site, aplicativo, redes sociais, antes de ir à loja?

\textbf{Luana:} Sigo bastante conta de moda e tendência, mas isso é mais pesquisa acadêmica minha do que pesquisa de compra propriamente.

\textbf{Danilo:} O que você faz quando isso acontece, pede ajuda, procura sozinha por mais tempo, desiste?

\textbf{Luana:} Se eu não acho, eu simplesmente sigo olhando outras coisa, não é uma frustração grande pra mim.

\textbf{Danilo:} Isso te faz sair da loja de mãos vazias com que frequência?

\textbf{Luana:} Quase sempre, na verdade, porque eu vou mais pra olhar do que pra comprar.

\textbf{Danilo:} Se você soubesse, antes de sair de casa, que a loja tem o que você quer, isso mudaria alguma coisa em como você compra?

\textbf{Luana:} Não mudaria muito pra mim, porque raramente eu vou com um objetivo de compra definido.

\textbf{Danilo:} Voltando à sua última compra, em algum momento um vendedor te ajudou? Como foi isso?

\textbf{Luana:} Não, resolvi tudo sozinha, nem cheguei a interagir com ninguém da loja.

\textbf{Danilo:} Tem alguma parte da compra em que você prefere não falar com ninguém e resolver sozinha?

\textbf{Luana:} Praticamente tudo, eu realmente prefiro observar e decidir sozinha, sem interferência.

\textbf{Danilo:} O que um vendedor resolve para você que você não resolveria sozinha?

\textbf{Luana:} Sinceramente, quase nada, eu quase nunca preciso de ajuda de vendedor.

\textbf{Danilo:} Já teve uma situação em que a ajuda de um vendedor atrapalhou mais do que ajudou, ou o contrário?

\textbf{Luana:} Já tive vendedor vindo falar comigo assim que eu entrei, isso me incomoda porque eu gosto de observar em paz antes de qualquer interação.

\textbf{Danilo:} Isso muda dependendo do tipo de peça ou do valor?

\textbf{Luana:} Não muito, eu sou sempre assim, independente do valor.

\textbf{Danilo:} Isso mudou com o tempo, você busca mais ou menos ajuda hoje do que buscava antes?

\textbf{Luana:} Busco menos, desde que eu comecei a estudar moda eu me sinto mais segura pra decidir sozinha.

\textbf{Danilo:} Como foi o momento de pagar na sua última compra? Teve fila?

\textbf{Luana:} Não teve, mas mesmo se tivesse eu acho que não me incomodaria muito, não é algo que eu preste muita atenção.

\textbf{Danilo:} De modo geral, o que você acha do processo entre decidir que vai levar uma peça e efetivamente sair da loja com ela?

\textbf{Luana:} Nunca parei muito pra pensar nisso, sinceramente, porque eu compro tão pouco que isso não é uma dor pra mim.

\textbf{Danilo:} Já desistiu de comprar algo por causa da fila ou da demora no caixa?

\textbf{Luana:} Não me lembro de isso ter acontecido.

\textbf{Danilo:} O que acontece, na prática, entre você decidir vou levar isso e você sair da loja com a peça em mãos?

\textbf{Luana:} É bem direto pra mim, vou até o caixa, pago e saio, não é algo que eu penso muito.

\textbf{Danilo:} Isso acontece com frequência ou foi uma exceção?

\textbf{Luana:} Como eu compro raramente, não dá pra dizer que é frequente nem exceção, é só como sempre foi as poucas vezes que aconteceu.

\textbf{Danilo:} Você já usou algum tipo de autoatendimento, supermercado, farmácia, cinema, sem passar por um funcionário?

\textbf{Luana:} Uso bastante, eu sou bem à vontade com esse tipo de coisa.

\textbf{Danilo:} Você confiaria em pagar por uma peça de roupa pelo celular, sem precisar falar com ninguém na loja?

\textbf{Luana:} Confiaria tecnicamente, sim, eu uso tecnologia sem problema. Mas eu ia sentir falta do ritual de pagar olhando a peça, sei lá, é bobagem talvez, mas faz parte da experiência pra mim.

\textbf{Danilo:} Como você se sente ao sair de uma loja com uma compra, já teve alguma experiência desconfortável com alarme, checagem de nota ou verificação na saída?

\textbf{Luana:} Nunca tive problema, mas confesso que eu já me senti desconfortável vendo outras pessoas passando por isso na minha frente, parece meio constrangedor.

\textbf{Danilo:} O que faria você confiar mais, ou menos, num processo assim?

\textbf{Luana:} Confiaria mais se fosse discreto, sem chamar atenção de quem tá por perto.

\textbf{Danilo:} Tem alguma coisa sobre comprar roupa em loja física que a gente não conversou e que você acha importante?

\textbf{Luana:} Acho que vale falar que hoje eu penso muito em sustentabilidade, prefiro brechó, e às vezes eu fico incomodada com o volume de roupa nova que essas loja grande produz. Sei que não é bem o foco da nossa conversa, mas é algo que pesa pra mim.

\textbf{Danilo:} Se você pudesse mudar uma única coisa em como as lojas funcionam hoje, o que seria?

\textbf{Luana:} Talvez trazer mais transparência sobre a origem das peça, mas eu entendo que isso foge um pouco do que vocês estão perguntando.

\textbf{Danilo:} Já viu, em outra loja, outro país ou outro tipo de serviço, alguma coisa que resolvesse um problema parecido com os que você mencionou?

\textbf{Luana:} Vi loja lá fora que tinha um espelho digital que mostrava outras combinação de look sem precisar trocar de roupa toda hora, achei bem inteligente.

\textbf{Danilo:} A Riachuelo está estudando uma etiqueta de segurança na peça que se comunica com o aplicativo da loja. Ela sabe se a peça foi paga e, quando o pagamento é confirmado pelo aplicativo, ela se destrava sozinha, sem precisar passar pelo caixa. O que chama sua atenção nisso, para o bem ou para o mal?

\textbf{Luana:} Achei tecnicamente interessante, mas sinceramente não é algo que resolve alguma dor que eu tenha, porque eu quase não compro e quando compro não é o pagamento que me incomoda.

\textbf{Danilo:} Em que situação você usaria isso, e em que situação não usaria?

\textbf{Luana:} Talvez eu usasse só por curiosidade, pra testar como funciona, mais do que por necessidade real.

\textbf{Danilo:} Isso resolve algum dos problemas que você mencionou antes, ou resolve outra coisa?

\textbf{Luana:} Não resolve nada do que eu mencionei, resolve um problema que eu acho que é mais de outro tipo de cliente, não o meu caso.

\textbf{Danilo:} O que precisaria existir, ou o que teria que te garantir, para você confiar nesse processo?

\textbf{Luana:} Não sei se eu chegaria a me importar o suficiente pra pensar nisso a fundo, sinceramente.

\textbf{Danilo:} Você imagina algum problema com uma peça saber que foi manuseada ou onde está dentro da loja?

\textbf{Luana:} Isso me deixa um pouco desconfortável, eu penso em quanto dado estão coletando sobre isso, mas eu não tenho uma opinião muito formada, é mais uma sensação.

